%% The Elicit installation manual — umbrella UC-001, "Consult the Installation Manual".
%%
%% Build it with docs/manual/build-manual.sh, which renders it through a TeX Live container
%% and stamps it with the platform version and build date (UC-002 BR-002). See README.md.
%%
%% Scope is provisioning a deployment (umbrella C-004), including the reports a site adds to
%% its surveys (UC-020, UC-021). Operating the console is the administrator's manual (Admin
%% UC-029). This document points at it and stops.

\documentclass[11pt, a4paper]{article}

\usepackage[T1]{fontenc}
\usepackage[a4paper, margin = 25mm, bottom = 28mm]{geometry}

\newcommand{\elicitProduct}{Elicit Platform}
\usepackage{elicit-brand}
\usepackage[hidelinks]{hyperref}

\hypersetup{
  pdftitle = {Elicit — Installing the Platform},
  pdfauthor = {Elicit Software},
  pdfsubject = {Installing, configuring and customizing the Elicit platform},
  colorlinks = true,
  linkcolor = elicitInkDark,
  urlcolor = elicitInkDark,
}

\graphicspath{{./}}

\begin{document}

%% ------------------------------------------------------------------ title page
\begin{titlepage}
\thispagestyle{empty}
\centering
\vspace*{3.5cm}

\includegraphics[width = 34mm]{images/elicit-logo.png}

\vspace{1.4cm}

{\color{elicitInkDarkest}\Huge\bfseries Installing the Platform\par}

\vspace{0.5cm}

{\color{elicitSlate}\Large The Elicit installation manual\par}

\vspace{0.4cm}

{\color{elicitSlate}\large Survey \textperiodcentered\ Admin \textperiodcentered\ FHHS \textperiodcentered\ Pedigree\par}

\vspace{2.2cm}

{\color{elicitSlate}\large
Version \textbf{\elicitVersion}\\[0.35em]
Built \textbf{\elicitBuildDate}\par}

\vfill

{\color{elicitSlate}\small
This manual describes the platform version named above. It is a release artifact of the\\
Elicit umbrella repository and is read before any Elicit service is running.\\[0.6em]
Elicit Software \textperiodcentered\ PolyForm Noncommercial License 1.0.0\par}

\vspace{1cm}
\end{titlepage}

%% ------------------------------------------------------------------ contents
\pagenumbering{roman}
\setcounter{page}{1}
\begin{elicitnote}[Who this manual is for]
Someone comfortable with PostgreSQL, a container runtime and an OIDC provider, who has not
seen Elicit before. Parts~I and II take a site from an empty machine to its first survey.
Part~III makes it the site's own: its brand, its languages and its reports, the last of which
is written for whoever builds a report service. Part~IV checks the result, and the appendix is
the reference for every setting.

Running the console day to day is the administrator's manual. This one stops where that
begins.
\end{elicitnote}

\vspace{1.2em}
{\color{elicitInkDarkest}\section*{Contents}}
\vspace{-0.6em}
\begingroup
\hypersetup{linkcolor = elicitCharcoal}
\makeatletter\@starttoc{toc}\makeatother
\endgroup

\clearpage
\pagenumbering{arabic}
\setcounter{page}{1}

\elicitpart{Getting oriented}{What an Elicit site is made of, how its pieces depend on each
other, and what has to exist before you install any of them.}

%% ================================================================== 1
\section{What Elicit is made of}
\label{sec:topology}

Elicit is a modular survey platform: a handful of applications that share one database, one
identity provider and one brand directory. A \term{site} is one installation of them. This
chapter names the pieces and how they connect, so that you can decide which of them your site
runs before you install anything.

\begin{elicitdiagram}{A site. Respondents reach Survey and administrators reach Admin. The
optional report services are called by Survey and Admin, never by a browser.}
  \node[elicit person, minimum width = 2.6cm] (resp) at (0, 0) {Respondent};
  \node[elicit person, minimum width = 2.6cm] (adm) at (6.4, 0) {Administrator};
  \node[elicit app, minimum width = 2.6cm] (survey) at (0, -1.8) {Survey};
  \node[elicit app, minimum width = 2.6cm] (admin) at (6.4, -1.8) {Admin};
  \node[elicit service, minimum width = 2.3cm] (idp) at (10.0, -1.25) {Identity provider};
  \node[elicit service, minimum width = 2.3cm] (smtp) at (10.0, -2.35) {SMTP relay};
  \node[elicit service, minimum width = 4.4cm] (db) at (3.2, -3.6)
    {PostgreSQL\\[-1pt]{\scriptsize the \val{survey} schema and the reporting schemas}};
  \node[elicit app, minimum width = 2.6cm] (fhhs) at (0, -5.4)
    {FHHS\\[-1pt]{\scriptsize report service}};
  \node[elicit app, minimum width = 2.6cm] (ped) at (3.6, -5.4)
    {Pedigree\\[-1pt]{\scriptsize draws pedigrees}};
  \begin{scope}[on background layer]
    \node[elicit group, fit = (fhhs) (ped), inner ysep = 12pt] (opt) {};
  \end{scope}
  \node[elicit group label, anchor = south west] at (opt.north west) {optional};
  \draw[elicit flow] (resp) -- node[elicit label] {access code} (survey);
  \draw[elicit flow] (adm) -- node[elicit label] {signs in} (admin);
  \draw[elicit flow] (survey) -- (db);
  \draw[elicit flow] (admin) -- (db);
  \draw[elicit flow] (admin.east) -- (idp.west);
  \draw[elicit flow] (admin.east) -- (smtp.west);
  \draw[elicit flow] ([xshift = -0.9cm]survey.south) -- node[elicit label, pos = 0.38]
    {reports} ([xshift = -0.9cm]fhhs.north);
  \draw[elicit flow] (fhhs) -- (db);
  \draw[elicit flow] (fhhs) -- (ped);
\end{elicitdiagram}

\subsection{The modules}

\begin{plaintable}{tab:modules}{The modules, what each one is, and what it needs}{@{}>{\raggedright\arraybackslash}p{0.16\linewidth} >{\raggedright\arraybackslash}p{0.40\linewidth} >{\raggedright\arraybackslash}p{0.37\linewidth}@{}}{\thead{Module} & \thead{What it is} & \thead{What it needs}}
\ui{Survey} & The subject-facing application. A respondent enters an \term{access code} and
answers the survey. It owns the shared \val{survey} schema and creates it. &
The database. Required. \\
\addlinespace
\ui{Admin} & The administration console: departments, users, subjects, access codes,
message templates, reports, and applying survey definitions. &
The database, the identity provider, an SMTP relay. Required. \\
\addlinespace
\ui{FHHS} & The report service of the Family Health History Survey, and the reference
example of a custom report (Chapter~\ref{sec:fhhs}). Optional. &
The database, Pedigree. \\
\addlinespace
\ui{Pedigree} & Draws the family-tree diagrams FHHS asks for. Optional. & Nothing. \\
\addlinespace
\ui{Author} & The authoring tool: designs surveys and exports the definition files Admin
applies. Optional, and installed from its own manual. &
A database of its own, the identity provider, a preview Survey. \\
\end{plaintable}

\begin{elicitnote}[Optional means optional]
A site that does not run the Family Health History Survey runs neither \ui{FHHS} nor
\ui{Pedigree}, and nothing outside Chapter~\ref{sec:fhhs} requires them. A site that only
delivers surveys, applying definition files written elsewhere, runs no \ui{Author}; the
authoring tool has an installation manual of its own.
\end{elicitnote}

\subsection{What they share}

\begin{itemize}
  \item \textbf{One PostgreSQL cluster.} A site's survey database carries the
        \val{survey} schema, the common reporting schema \val{surveyreport}, and one
        reporting schema per installed survey, which Survey creates itself. Survey, Admin
        and FHHS all connect to it.
  \item \textbf{One identity provider.} Admin authenticates against its own OIDC client.
        Survey authenticates no one: a respondent enters an access code, which is not a
        credential in the identity provider.
  \item \textbf{One brand directory,} mounted read-only into Survey and Admin.
\end{itemize}

Translations are optional. Survey and Admin carry English in the image. Latin American Spanish
(\val{es-419}) and Arabic (\val{ar}) come in a curated translations directory that
ElicitSoftware maintains beside each release, which a site mounts read-only, along with any
files of its own. A site that mounts nothing offers English only (Chapter~\ref{sec:i18n}).

\subsection{The order of the first start}

Survey creates and migrates the shared \val{survey} schema. Admin adds its own tables to that
schema, and FHHS adds its grants and indexes on top of both. So on a \emph{new} database:

\begin{enumerate}
  \item \ui{Survey} starts and becomes ready. Readiness is what tells you its migrations
        have finished.
  \item \ui{Admin} starts and becomes ready. Its migrations run against the schema Survey
        created.
  \item \ui{FHHS} starts. Started beside Admin rather than after it, the two modules' first
        migrations grant on the same tables at the same moment, and one of them fails.
  \item \ui{Pedigree} starts whenever you like: it has no database.
\end{enumerate}

\begin{elicitnote}[The order constrains the first start only]
It exists because of schema creation. Once the schema is migrated the modules may be started
in any order and may restart independently. The compose file enforces the order for you;
without it, the order is yours to keep (Chapter~\ref{sec:nocompose}).
\end{elicitnote}

\subsection{Ports}

Every application listens on port \val{8080} \emph{inside} its container. The host ports
below are what the shipped compose file publishes.

\begin{plaintable}{tab:ports}{Published host ports in the shipped compose file}{@{}>{\raggedright\arraybackslash}p{0.13\linewidth} >{\raggedright\arraybackslash}p{0.28\linewidth} >{\raggedright\arraybackslash}p{0.49\linewidth}@{}}{\thead{Host port} & \thead{Service} & \thead{Reached by}}
\val{8080} & Survey & Respondents, in a browser. \\
\val{8081} & Admin & Administrators, in a browser. \\
\val{8082} & FHHS & Survey and Admin, for reports. Published for local testing only. \\
\val{8083} & Pedigree & FHHS. Published for local testing only. \\
\val{8084} & Author & Survey authors, in a browser. See the Author installation manual. \\
\val{8085} & Author's preview Survey & Survey authors, in a browser. \\
\val{5452} & PostgreSQL (\val{5432} inside) & The applications, and you. \\
\val{8180} & Keycloak (bundled) & Browsers \emph{and} the application containers. \\
\val{8025} & Mailpit web interface (bundled) & You. \\
\val{1025} & Mailpit SMTP (bundled) & Admin. \\
\val{16686} & Jaeger user interface (bundled) & You. \\
\val{4317} & Jaeger OTLP/gRPC collector (bundled) & Every application. \\
\end{plaintable}

\begin{elicitcaution}[Health checks address the container, not the host]
Each service's health check targets \val{http://localhost:8080/q/health/ready} \emph{inside}
the container. That is the same address for every application, whichever host port it is
published on. A health check pointed at the published host port is the most common
self-inflicted installation failure.
\end{elicitcaution}

\subsection{Metrics and traces}

Every application exposes Prometheus-format metrics at \opt{/q/metrics} and health at
\opt{/q/health/live} and \opt{/q/health/ready}, and exports OpenTelemetry traces to a
collector.

\begin{elicitnote}[There is no metrics stack in this repository]
Nothing in the shipped compose file scrapes \opt{/q/metrics}, and no Prometheus service runs.
The endpoint is there for your site's own scraper. The bundled Jaeger collects traces only, so
its \ui{Monitor} tab, which would read metrics, is inert.
\end{elicitnote}

%% ================================================================== 2
\section{Before you start}
\label{sec:before}

Check each of these for the modules your site will run. A missing prerequisite is far cheaper
to find now than halfway through an installation.

\begin{description}
  \item[A PostgreSQL 17 cluster] in which you can create roles, databases and schemas.
  \item[A container runtime,] with compose for the bundled path (Chapter~\ref{sec:compose}),
        or the container platform your site already runs (Chapter~\ref{sec:nocompose}).
        Elicit is released as container images only.
  \item[An OIDC provider] in which you can create a confidential client and three roles,
        reachable at \emph{one} address from both the browser and the application containers.
  \item[An SMTP relay] for Admin's invitations and reminders. Admin is the only module that
        sends mail.
  \item[A directory for the brand,] readable by every module that mounts it, and the same for
        translations if the site offers any language besides English (Chapter~\ref{sec:i18n}).
  \item[Free host ports] for the modules you will run (Table~\ref{tab:ports}).
\end{description}

\begin{elicitcaution}[One address for the identity provider]
The issuer address you configure must resolve to the same provider \emph{from the browser}
and \emph{from inside the container}. When it does not, sign-in fails in a way that looks
like an application fault: the browser is redirected somewhere it cannot reach, or, on a host
that maps the provider's name to itself, lands on whichever service listens on that port. The
bundled example uses \val{host.docker.internal}, which Docker Desktop resolves to the host
from inside a container and to the loopback address on the host, so one address works from
both.
\end{elicitcaution}

The repository bundles a database, an identity provider, a mail catcher, an SFTP server and a
trace collector so that an evaluation site works in one command. Each ships with a documented
password; read Chapter~\ref{sec:secrets} before any of them reaches a real site.

\elicitpart{Installing}{From an empty database cluster to a site with its first survey: the
database, the identity provider, the secrets, the two ways to start the applications, the
first sign-in, and applying a survey definition.}

%% ================================================================== 3
\section{The database}
\label{sec:database}

Elicit creates no database roles: every role must exist before any module starts. Nor do you
create any tables. You create two empty schemas, and everything inside them comes from the
modules' migrations on their first start.

\subsection{The roles}

\begin{plaintable}{tab:roles}{The database roles Elicit expects}{@{}>{\raggedright\arraybackslash}p{0.24\linewidth} >{\raggedright\arraybackslash}p{0.14\linewidth} >{\raggedright\arraybackslash}p{0.56\linewidth}@{}}{\thead{Role} & \thead{Kind} & \thead{Used by}}
\val{dbowner} & group & Holds the create privileges the schema owner inherits. \\
\val{dbuser} & group & Holds the connect and usage privileges every application role
inherits. \\
\addlinespace
\val{elicit_owner} & login & The \emph{schema owner}. Every module's Flyway migrations run as
this role. \\
\val{survey_user} & login & The application connection of Survey and of FHHS. \\
\val{surveyadmin_user} & login & Admin's application connection. \\
\val{surveyreport_user} & login & Reads the reporting schemas, for a report service of the
site's own (Chapter~\ref{sec:reports}). \\
\end{plaintable}

Run this once for the cluster, substituting a password of your own for each placeholder.

\begin{elicitsql}
CREATE ROLE dbowner WITH NOLOGIN SUPERUSER INHERIT
  NOCREATEDB NOCREATEROLE NOREPLICATION;

CREATE ROLE dbuser WITH NOLOGIN NOSUPERUSER INHERIT
  NOCREATEDB NOCREATEROLE NOREPLICATION;

CREATE ROLE elicit_owner WITH LOGIN SUPERUSER INHERIT
  NOCREATEDB NOCREATEROLE NOREPLICATION
  ENCRYPTED PASSWORD '<owner password>';

CREATE ROLE survey_user WITH LOGIN NOSUPERUSER INHERIT
  NOCREATEDB NOCREATEROLE NOREPLICATION
  ENCRYPTED PASSWORD '<survey password>';
GRANT dbuser TO survey_user;

CREATE ROLE surveyadmin_user WITH LOGIN NOSUPERUSER INHERIT
  NOCREATEDB NOCREATEROLE NOREPLICATION
  ENCRYPTED PASSWORD '<admin password>';
GRANT dbuser TO surveyadmin_user;

CREATE ROLE surveyreport_user WITH LOGIN NOSUPERUSER INHERIT
  NOCREATEDB NOCREATEROLE NOREPLICATION
  ENCRYPTED PASSWORD '<report password>';
GRANT dbuser TO surveyreport_user;
\end{elicitsql}

\begin{elicitcaution}[The schema owner is a superuser here]
The role script creates \val{elicit_owner} as a superuser, because that is what the modules'
migrations are tested with. A site whose policy forbids a superuser login role can narrow it,
but must then verify the narrowed role against its release. The role needs to:
\begin{itemize}
  \item create and alter objects in the \val{survey} and \val{surveyreport} schemas;
  \item create a schema for each installed survey, which is why \val{CREATE ON DATABASE} below
        is required rather than a courtesy;
  \item own and alter what it creates there.
\end{itemize}
\end{elicitcaution}

\subsection{The site's survey database}

Create a database named \val{survey}, then run this inside it. \val{surveyreport} is the
common reporting schema: it holds only the date and status dimensions every survey shares.
Each survey's own star schema, named \val{report_...} after the survey and recorded on
\val{survey.surveys.report_schema}, is created by Survey when the survey is first built, not
here; the \val{CREATE ON DATABASE} grant is what allows it.

\begin{elicitsql}
GRANT CONNECT, CREATE ON DATABASE survey TO dbowner;
GRANT CONNECT ON DATABASE survey TO dbuser;

GRANT USAGE, CREATE ON SCHEMA public TO dbowner;
GRANT USAGE ON SCHEMA public TO dbuser;

CREATE SCHEMA survey AUTHORIZATION elicit_owner;
GRANT USAGE ON SCHEMA survey TO survey_user;
GRANT USAGE ON SCHEMA survey TO surveyadmin_user;
GRANT USAGE ON SCHEMA survey TO surveyreport_user;

CREATE SCHEMA surveyreport AUTHORIZATION elicit_owner;
GRANT USAGE ON SCHEMA surveyreport TO survey_user;
GRANT USAGE ON SCHEMA surveyreport TO surveyadmin_user;
GRANT USAGE ON SCHEMA surveyreport TO surveyreport_user;

GRANT USAGE, CREATE ON SCHEMA surveyreport TO dbowner;
ALTER DEFAULT PRIVILEGES IN SCHEMA surveyreport
  GRANT SELECT, INSERT, UPDATE, DELETE ON TABLES TO dbowner;
ALTER DEFAULT PRIVILEGES IN SCHEMA surveyreport
  GRANT USAGE ON SEQUENCES TO dbowner;
\end{elicitsql}

Both schemas are now empty. That is correct: every table, sequence, index and grant inside
them is created by a module's migrations the first time it starts.

\begin{elicitnote}[The bundled database image does all of this for you]
\file{elicitsoftware/elicit_db} runs the role script and the schema script when its data
directory is first initialized. An initialization script runs \emph{only} against an empty
data directory, so a cluster that already exists is yours to prepare with the two scripts
above. The image's scripts carry the documented development passwords listed in
Chapter~\ref{sec:secrets}.
\end{elicitnote}

%% ================================================================== 4
\section{The identity provider}
\label{sec:oidc}

Admin authenticates against a confidential OIDC client. Survey has no client: a respondent
enters an access code, which is issued by Admin and is not a credential in your identity
provider.

\subsection{The client}

\begin{plaintable}{tab:clients}{The OIDC client}{@{}>{\raggedright\arraybackslash}p{0.22\linewidth} >{\raggedright\arraybackslash}p{0.20\linewidth} >{\raggedright\arraybackslash}p{0.52\linewidth}@{}}{\thead{Client} & \thead{Type} & \thead{Redirect}}
\val{elicit-admin} & confidential & The console's address, for example
\val{http://localhost:8081/*}. \\
\end{plaintable}

The client needs a secret, which reaches the application as \opt{OIDC_CLIENT_SECRET}. Admin
accepts browser sign-ins and bearer tokens on the same client.

\screenshot{01-keycloak-clients}{The bundled realm's clients. A site creates the console's
client in its own provider, under whatever name it configures Admin with.}

\screenshot{02-keycloak-client}{The console's client. The redirect URI must cover the address
a browser reaches the console at, which is not necessarily the address the container knows.}

\subsection{The roles}

\begin{plaintable}{tab:roles-oidc}{The roles the applications check}{@{}>{\raggedright\arraybackslash}p{0.20\linewidth} >{\raggedright\arraybackslash}p{0.18\linewidth} >{\raggedright\arraybackslash}p{0.56\linewidth}@{}}{\thead{Role} & \thead{On client} & \thead{What it permits}}
\val{elicit_admin} & \val{elicit-admin} & The console's administrative screens. \\
\addlinespace
\val{elicit_user} & \val{elicit-admin} & The console's day-to-day screens, without the
administrative ones. \\
\addlinespace
\val{elicit_importer} & \val{elicit-admin} & The subject-import endpoints of the
integration API, for a machine client presenting a bearer token. \\
\end{plaintable}

At least one account must hold \val{elicit_admin} before the first sign-in
(Chapter~\ref{sec:firstsignin}), or nobody can create the first department.

\screenshot{03-keycloak-roles}{The roles on the console's client. A site creates the same
ones in its own provider.}

\subsection{The bundled example}

The repository ships a Keycloak realm at \file{keycloak/elicit-realm.json}, carrying the
console's client and its roles, the integration client \val{import-service}, Author's client,
and the accounts \val{admin}, \val{user} and \val{author}. It is a worked example and an
evaluation provider, not a production one: its client secret is the literal string
\val{secret}, and its administrator account is \val{admin}/\val{admin}.

\subsection{Sourcing roles from the database instead}

A site that cannot manage Elicit's roles in its identity provider can set
\opt{elicit.authorization.mode} to \val{DATABASE} on Admin. That makes the console consult
its own \val{survey.user_roles} table as a fallback, and shows the role-assignment section of
the user editor. It does \emph{not} turn OIDC off: sign-in is unchanged, and roles the provider
supplies are still honored. In this mode the identity provider's account names must match the
console's user records.

%% ================================================================== 5
\section{Replacing the development secrets}
\label{sec:secrets}

The repository and its images ship with working credentials so that an evaluation site starts
in one command. Every one of them has to be replaced before a real site goes live.

\begin{elicitsecurity}[Every one of these is published]
Each value below is in this repository and in this manual. A site that leaves one as shipped
is running with a credential anyone can read.
\end{elicitsecurity}

\begin{plaintable}{tab:secrets}{Development credentials that must be replaced}{@{}>{\raggedright\arraybackslash}p{0.29\linewidth} >{\raggedright\arraybackslash}p{0.21\linewidth} >{\raggedright\arraybackslash}p{0.43\linewidth}@{}}{\thead{What} & \thead{Shipped value} & \thead{Where it comes from}}
Schema owner password & \val{SURVEYPW} & The database image's role script;
\opt{ELICIT_OWNER_DB_PASSWORD} in the compose file, for every module. \\
\addlinespace
Survey's and FHHS's connection password & \val{SURVEYPW} & The database image's role script;
\opt{ELICIT_DB_PASSWORD} on the \val{survey} and \val{fhhs} services in the compose file. \\
\addlinespace
Admin's connection password & \val{SAOWNERPW} & The database image's role script;
\opt{ELICIT_DB_PASSWORD} in the compose file. \\
\addlinespace
Reporting password & \val{SRUSERPW} & The database image's role script. No shipped module
presents it; a report service of the site's own does (Chapter~\ref{sec:reports}). \\
\addlinespace
OIDC client secret, the console & \val{secret} & \opt{OIDC_CLIENT_SECRET} in the compose
file, and the \val{elicit-admin} client of the bundled Keycloak realm. \\
\addlinespace
OIDC client secret, integration client & \val{importer} & The \val{import-service} client of
the bundled Keycloak realm, used by a machine client presenting a bearer token. \\
\addlinespace
Database superuser & \val{survey}/\val{admin} & The bundled database service. \\
\addlinespace
Keycloak administrator & \val{admin}/\val{admin} & The bundled identity provider. \\
\addlinespace
SFTP account & \val{fhhs_user}/\val{pass} & The bundled SFTP service, for FHHS report
delivery. \\
\end{plaintable}

\begin{elicitcaution}[Every module needs both variables]
Survey, Admin and FHHS read their database passwords from \opt{ELICIT_DB_PASSWORD} and
\opt{ELICIT_OWNER_DB_PASSWORD}, and none of them ships a value, so no module ever connects with
a published password by default. A module started without them does not stop at once, though.
The database refuses it, and the module logs \val{no password was provided}, reports itself
not ready, and retries \opt{quarkus.flyway.owner.connect-retries} times with waits that double
up to two minutes. Only then does it exit with \val{Failed to start application}: about eight
minutes after it started, at the shipped ten retries. A module that sits unhealthy with that
message in its log is missing these variables.

Before 3.0.0, Survey and FHHS carried the literal \val{SURVEYPW}, so a site upgrading from an
earlier release must now set both variables on them, or \opt{quarkus.datasource.password} and
\opt{quarkus.datasource.owner.password}, which take precedence.
\end{elicitcaution}

\begin{elicitnote}[The two Admin variables read backwards, and are correct]
On the \val{admin} service the compose file sets \opt{ELICIT_DB_PASSWORD}\val{=SAOWNERPW} and
\opt{ELICIT_OWNER_DB_PASSWORD}\val{=SURVEYPW}. That is right: Admin's application connection
is \val{surveyadmin_user}, whose password is \val{SAOWNERPW}, and its migrating connection is
\val{elicit_owner}, whose password is \val{SURVEYPW}. The names simply do not read that way.
Do not swap them.
\end{elicitnote}

Two more things a new installation ships with are not secrets but are just as public: the
seeded accounts \val{admin} and \val{user}, which Admin warns about on every screen until
they are renamed (Chapter~\ref{sec:firstsignin}).

%% ================================================================== 6
\section{Installing with Docker Compose}
\label{sec:compose}

The umbrella repository carries \file{docker-compose.yml}, the default brand and the curated
translations directory, \file{elicit-i18n/}. Together they are both a working site and a
worked example of one. The compose file mounts the brand and the translations read-only into
Survey and Admin (Chapters~\ref{sec:brand} and~\ref{sec:i18n}).

\subsection{What the file defines}

Six application services and five supporting services bundled for evaluation:

\begin{itemize}
  \item the site: \val{survey}, \val{admin}, and the Family History reports \val{fhhs} and
        \val{pedigree} (Chapter~\ref{sec:fhhs});
  \item the authoring tool: \val{author} and its preview, \val{author-survey}, which the
        Author installation manual covers;
  \item for evaluation: \val{db}, \val{keycloak}, \val{mailpit}, \val{sftpServer} and
        \val{jaeger}.
\end{itemize}

Every application image is \file{elicitsoftware/<name>} at one tag,
\val{${ELICIT_TAG:-latest}}, so the applications always run as a matched set from one
release. The tag defaults to \val{latest}; you run a published release by naming its version
(Section~\ref{sec:compose-up}). The database image, \file{elicitsoftware/elicit_db}, is
versioned on its own and is not on the tag.

\subsection{Before bringing it up}

\begin{enumerate}
  \item Replace the development secrets (Chapter~\ref{sec:secrets}).
  \item Set \opt{accessCode.autoRegister} to \val{false} on \val{survey}. The file ships it
        \val{true} so that an evaluator can type any access code; on a site that delivers real
        surveys, that would let anyone in.
  \item Point the brand volume at your own brand directory, if you have one
        (Chapter~\ref{sec:brand}).
  \item Point the translations volume at your own directory if you have one, or remove it for
        an English-only site, and set \opt{i18n.offered.locales} on any application that should
        offer fewer languages than are mounted (Chapter~\ref{sec:i18n}).
  \item Remove any bundled service your site replaces with its own, and point the
        applications at yours using the settings in Appendix~\ref{sec:config}.
  \item Remove \val{fhhs} and \val{pedigree}, and Admin's dependency on \val{pedigree}, if
        your site does not run the Family Health History Survey. Remove \val{author} and
        \val{author-survey} if it authors no surveys.
\end{enumerate}

\subsection{Bringing it up}
\label{sec:compose-up}

\begin{elicitcmd}
ELICIT_TAG=3.0.0-beta.1 docker compose up -d
./status.sh
\end{elicitcmd}

Name the release you are installing in \val{ELICIT_TAG}, either on the command line as above
or once, in a \file{.env} file beside the compose file:

\begin{elicitcmd}
ELICIT_TAG=3.0.0-beta.1
\end{elicitcmd}

A pre-release (\val{3.0.0-beta.1}, \val{3.0.0-rc.1}) is published under its exact version
only; \val{latest} and the \val{3.0} tag move only when a final version ships. Left unset, the
tag is \val{latest}, which is what a local build produces.

That is the whole procedure. The database starts and becomes healthy, Survey creates and
migrates the shared schema, Admin migrates on top of it, and the remaining services start as
their dependencies become healthy.

\begin{elicitnote}[No restart sequence is required]
Older instructions described bringing the stack up and then restarting it twice. That no
longer applies. Survey creates the schema, a survey definition is applied through the console
afterwards (Chapter~\ref{sec:definition}), and FHHS becomes ready by itself on its next
health probe. FHHS waits for Admin, not the other way round, so the console comes up on a
new database and the definition can be applied through it while FHHS is still waiting for
its survey.
\end{elicitnote}

If a published host port is taken, change it, and change the matching redirect URI in your
identity provider too (Chapter~\ref{sec:oidc}).

\subsection{Resetting}

To start again from an empty database, run

\begin{elicitcmd}
./resetDatabase.sh V3
\end{elicitcmd}

which stops the stack and deletes the database's data directory. Everything in the database is
lost, so use it only on an evaluation site.

%% ================================================================== 7
\section{Installing without Compose}
\label{sec:nocompose}

Elicit is released as container images only, \file{elicitsoftware/<name>} at the release's
tag; no runnable Java artifact is distributed. If your site runs its containers under something
other than the bundled compose file, such as Podman, Kubernetes or a cloud container service,
start the same images with the procedure below, once per module, against your own database
(Chapter~\ref{sec:database}) and identity provider (Chapter~\ref{sec:oidc}).

\begin{enumerate}
  \item Pull the module's image at the tag of the release you are installing. Every module
        runs at the same tag.
  \item Write its configuration (datasource, identity, mail, the brand path, and anything else
        your site changes) as environment variables, by the rules in
        Appendix~\ref{sec:config}, and mount the brand and translations directories read-only.
  \item Run it under your container platform, publishing its container port \val{8080}.
  \item Start \ui{Survey} first and wait for \opt{/q/health/ready} to report healthy. That is
        what confirms its migrations have completed.
  \item Start \ui{Admin} and wait for its readiness endpoint.
  \item Start \ui{FHHS}, if your site runs it, only once Admin is ready; start \ui{Pedigree}
        whenever you like (Chapter~\ref{sec:fhhs}).
\end{enumerate}

\begin{elicitcaution}[Nothing sequences this for you]
Unlike the compose file, your platform does not know the start order. A module started before
Survey has migrated a \emph{new} database fails its own migration on an object that does not
exist yet, and does not start. Nothing is corrupted: let Survey finish, then start the module
again.
\end{elicitcaution}

\subsection{Running more than one instance}

Instances of the same module must agree on their datasource, their brand path and their
translations: mount the same directories and set the same \opt{i18n.offered.locales} on each.
Instances of Survey that disagreed, behind one load balancer, would offer a different set of
languages depending on which one answered. Keep the proxy-forwarding settings on
(Table~\ref{tab:quarkus-http}), or sign-in redirects are built with the container's address
instead of the site's.

%% ================================================================== 8
\section{The first sign-in}
\label{sec:firstsignin}

A new installation seeds no survey and no department. It does seed two accounts, \val{admin}
and \val{user}, so that the console can be signed into at all.

\screenshot{04-sign-in}{Signing in. The console never asks for a password itself: it sends you
to the identity provider and takes the answer back.}

\screenshot{05-first-sign-in}{What a new installation looks like on the first sign-in: the
blocking dialog, the notice that no survey is installed, and the warning about the two seeded
accounts. All three clear during this chapter and the next.}

\begin{enumerate}
  \item Sign in to Admin with an account holding \val{elicit_admin}.
  \item A dialog blocks the console until a department exists. Follow its link to
        \ui{Departments} and create one; it is assigned to you as its creator.
  \item Rename the two seeded accounts to real people. Until both are renamed, Admin logs a
        warning at every start and shows a banner on every console screen. Under
        \opt{elicit.authorization.mode}\val{=DATABASE}, rename the matching identity-provider
        accounts to the same usernames.
  \item Write at least one message template, putting \val{<ACCESS_CODE>} wherever a subject's
        access code belongs, for example in the login link
        \val{https://<your survey host>/login/<ACCESS_CODE>}, which opens Survey with the
        code filled in.
\end{enumerate}

\screenshot{06-departments-empty}{Departments, before there are any. The dialog's link brings
you here.}

\screenshot{07-new-department}{The first department. \ui{Default message ID} is prefilled with
\val{1}, which on a new installation names a template that does not exist yet: write the
template first, or correct this field once you have.}

\screenshot{08-console-ready}{The console with a department. The dialog is gone; the two
banners stay until a survey is applied and the seeded accounts are renamed.}

\begin{elicitnote}[An account with no department can only log out]
That is deliberate. A user holding \val{elicit_user} but belonging to no department has
nothing to do in the console until an administrator assigns them one.
\end{elicitnote}

\begin{elicitcaution}[A department created in SQL is assigned to nobody]
Creating a department with a direct \val{INSERT} skips the assignment the console makes.
Assign it to its users separately, or create it through \ui{Departments}.
\end{elicitcaution}

%% ================================================================== 9
\section{Installing a survey definition}
\label{sec:definition}

A new site has no survey; surveys arrive as definition files, exported from Author and applied
in Admin.

\begin{enumerate}
  \item In Admin, open \ui{Apply Survey Definition}.
  \item Upload a \file{.elicit} file. For the Family Health History Survey that is
        \file{FHHS/family-history-survey.elicit}.
  \item The console installs the survey the first time and updates it in place afterwards,
        matching elements by their keys.
  \item As part of the apply, the console asks Survey to build the reporting star schema
        of the survey just applied. Survey creates a schema of the survey's own, named
        \val{report_} followed by the survey's name in lower case (for the Family Health
        History Survey, \val{report_family_history_survey}), records the name on
        \val{survey.surveys.report_schema}, and fills it with the dimension tables, the fact
        table and the views. \val{surveyreport} itself keeps only the shared date and status
        dimensions.
\end{enumerate}

\screenshot{09-apply-definition}{Apply Survey Definition. The file identifies itself, so the
same screen installs a survey and updates one.}

\screenshot{10-apply-result}{The result. It reports what was installed and, on the last line,
that the reporting schema was rebuilt, which is why no Survey restart is needed.}

\begin{elicitnote}[No Survey restart is needed]
The rebuild happens during the apply: Admin calls Survey, at \opt{elicit.survey.url}, unless
\opt{elicit.survey.etl-build.enabled} turns the call off. The outcome is reported in the apply
result and never fails the apply. A site that restarts Survey after every apply, or runs it with
\opt{elicit.etl.enabled}\val{=false}, can turn the call off.
\end{elicitnote}

To confirm the rebuild from the database:

\begin{elicitcmd}
docker exec elicit-db-1 psql -U survey -d survey \
  -c "select name, report_schema from survey.surveys;" \
  -c "select count(*) from report_family_history_survey.dim_step;"
\end{elicitcmd}

\val{report_schema} must be set and the count non-zero. The fact table stays empty until
respondents finish; \val{fact_respondents} is a view over \val{survey.respondents}, so it
follows every respondent as they progress.

The schema's name belongs to the site and is never carried in a definition file. An
administrator can rename it with \ui{Rename} on Admin's \ui{Export Survey Definition} screen.
Nothing inside Elicit needs changing after a rename, but queries outside Elicit that name the
old schema do.

If your site runs FHHS, it has been reporting not-ready and answering report requests with
\val{503} until now, naming the survey it is waiting for. It goes healthy by itself on its
next probe; nothing needs restarting (Chapter~\ref{sec:fhhs}).

\elicitpart{Customizing}{Making the site its own: its brand, the languages it offers, and the
reports its surveys end with. The reports chapters are written for whoever builds a report
service, and install the two that ship with Elicit.}

%% ================================================================== 10
\section{External branding}
\label{sec:brand}

Survey and Admin each carry the Elicit default brand inside the image. To give a site its own
colors, typefaces, logos and organization name, mount a brand directory and point
\opt{brand.file.system.path} at it. Nothing is rebuilt.

\subsection{The layout}

\begin{elicitfile}
<brand>/
  brand-config.json                name, organization, logo files
  brand-info.json                  description for page metadata
  theme.css                        imports the two stylesheets
  colors/brand-colors.css          --brand-* color properties
  typography/brand-typography.css  --brand-* typeface properties
  typography/typography.json       the typefaces, described
  images/HorizontalLogo.png        the logo in the header
  images/icon-white.png            the icon on a dark background
  images/favicon.ico               the browser tab's icon
\end{elicitfile}

Start from the repository's \file{elicit-brand/} and change only the files that differ.

\subsection{A brand may be partial}

Each file is looked for on its own: first in the mounted directory, then in a \file{brand}
directory beside the application, then in the default brand inside the image. A brand that
changes only the colors needs only \file{colors/brand-colors.css}, and a \file{theme.css} that
imports it.

A brand sets \val{--brand-*} custom properties and nothing else. Each application's own
stylesheet maps its theme's tokens onto them, which is what lets one brand directory serve both
applications unchanged.

\begin{elicitcaution}[A partial brand fails quietly]
A file the mount lacks, or names wrongly, is taken from the default brand without an error. The
page renders, in the default typeface or logo, and raises no error. The \ui{Branding} page
(below) is the place that shows it.
\end{elicitcaution}

\subsection{Mounting one}

Mount the directory read-only on Survey and on Admin, and set the path on each. In the compose
file:

\begin{elicitcmd}
volumes:
  - ./my-brand/:/opt/brand:ro
environment:
  brand.file.system.path: /opt/brand
\end{elicitcmd}

No application writes to the brand directory, so read-only costs nothing. The setting is per
application, so Survey and Admin may be given different brands if a site wants its
respondent-facing pages to look different from its console.

\screenshot{17-branded-header}{A branded application. The header logo, the colors and the
typefaces all come from the mounted directory.}

\subsection{Changing a brand on a running site}

Each application reads its brand once and holds it, so editing a file in the mount changes
nothing by itself. Restart the application to pick the change up. Admin's \ui{Branding} page
also has \ui{Reload brand}, which makes the console read its brand again without a restart; it
does not reach Survey, which is a separate service.

\subsection{Translated brand text}

The organization name in the header and the brand's description are translated \emph{in the
brand}, not in the applications' translation files. Add a \val{localized} block, keyed by
language tag, to \file{brand-config.json} for the name and the organization:

\begin{elicitfile}
{
  "name": "Elicit Default Brand",
  "organization": "Elicit Software",
  "localized": {
    "es-419": { "organization": "Elicit Software" },
    "ar":     { "organization": "Elicit Software" }
  }
}
\end{elicitfile}

A block of the same shape in \file{brand-info.json} translates the \val{description}. Each text
is looked up by the exact tag (\val{es-419}), then by the language alone (\val{es}), then falls
back to the untranslated value. Leave the base \val{name} as it is: the applications derive
the brand's technical key from it.

\subsection{Checking that it took}

Admin's \ui{System} section has a \ui{Branding} page. It names the brand directory that was
configured and whether it exists, and says, for each asset, whether it came from the mount, the
local directory or the default brand. It is the fastest way to find a mount that is one
directory level off.

\screenshot{13-system-branding}{The Branding page. It names the directory that resolved and,
for each asset, whether it came from the mount or from the default brand.}

%% ================================================================== 11
\section{Translations}
\label{sec:i18n}

Survey and Admin carry their own interface text, the \term{chrome}, in English inside the
image, and in no other language. Every other language is a file the site supplies, in a
translations directory mounted read-only. ElicitSoftware maintains a curated directory beside
each release, \file{elicit-i18n/} in the umbrella repository, which today holds Latin American
Spanish (\val{es-419}) and Arabic (\val{ar}) for both applications; the compose file mounts it.

An application offers English plus whatever is mounted for it, narrowed by
\opt{i18n.offered.locales} where that is set (Table~\ref{tab:common}). The language selector
appears in the header only when more than one language is offered, so a site that mounts
nothing never shows it.

\begin{elicitcaution}[Survey content is translated separately]
This mechanism translates the chrome: buttons, labels, page titles, notifications, validation
messages and PDF headers. The survey's own content is translated too, but not from here. That
content covers:
\begin{itemize}
  \item question text and answer options;
  \item step and section names;
  \item tooltips and validation messages.
\end{itemize}
Those translations are written for one survey in the authoring tool and arrive at a site inside
its definition file. A respondent is offered a content language only when the survey publishes
it \emph{and} the site offers that language for the chrome. So the mount and
\opt{i18n.offered.locales}, which govern the chrome, also decide which content languages can
reach a respondent. Message templates are outside both mechanisms and appear in
the language they were authored in. A report service is sent the respondent's language and
writes its report in it, or not, as it chooses (Chapter~\ref{sec:reports}).
\end{elicitcaution}

\subsection{The translations directory}

One directory serves both applications. Each reads only its own sub-directory,
\file{<i18n.file.system.path>/<i18n.app.name>/}: by default \file{/i18n/survey/} and
\file{/i18n/admin/}.

\begin{elicitfile}
elicit-i18n/                     mounted read-only at /opt/i18n
  i18n-config.json               text direction and size per language
  survey/
    translations_es_419.properties
    translations_ar.properties
  admin/
    translations_es_419.properties
    translations_ar.properties
\end{elicitfile}

Each \file{translations_<tag>.properties} file adds its language to that application. The
applications only read the directory, on the server: nothing in it is served to a browser, and
nothing is written to it.

The curated directory holds the wording a release was built and tested with, and is versioned
with it, so mount the copy that matches the version you run. A site may mount it whole, mount
only some of its files, or mount files of its own (Section~\ref{sec:commission}). Anything
beyond the curated files is the site's to maintain.

\screenshot{18-language-selector}{The selector, open on an application offering all three
languages of the curated directory. Each row is labeled in its own language. The selector appears
only when there is something to select, so a site offering only \val{en} never sees it.}

\subsection{How a text is resolved}

Text by text: the file for the exact tag (\val{es-419}), then the file for the language alone
(\val{es}), then English. A mounted file needs only the texts that differ from English; a text
it does not carry falls back to English rather than showing blank. A file holding one corrected
text is therefore a valid way to change one wording. A text missing everywhere shows as
\val{!key!} and is logged once.

\subsection{Choosing which languages to offer}
\label{sec:language}

\begin{enumerate}
  \item Mount the translations directory read-only on Survey and Admin, and set
        \opt{i18n.file.system.path} to the mount point. The compose file already mounts
        \file{elicit-i18n/} at \val{/opt/i18n}. Leave it unmounted for an English-only
        site.
  \item If an application should offer fewer languages than are mounted, set
        \opt{i18n.offered.locales} on it to a comma-separated list, for example
        \val{en,es-419}. Leave it unset to offer everything mounted.
  \item Restart each application. Both the directory and the list are read at startup.
\end{enumerate}

A language that is not offered is not offered anywhere: the selector hides it, a \opt{?lang=}
naming it is refused, and survey content is not served in it, because content follows the
chrome. \opt{i18n.offered.locales} cannot add a language: a tag it lists with no file behind
it is logged at startup and ignored.

\begin{elicitcaution}[The tag is punctuated with a hyphen; the file name with an underscore]
A language tag is written with a hyphen everywhere an operator sets one: \val{es-419},
\val{ar}. That covers \opt{i18n.offered.locales}, \optlit{i18n.direction.<tag>},
\optlit{i18n.font-scale.<tag>} and a \opt{?lang=} parameter. The file names in the
translations directory use an underscore instead (\file{translations_es_419.properties}). A
file named with a hyphen is not recognized as a language.
\end{elicitcaution}

The mount and \opt{i18n.offered.locales} are set per application, not per site. Survey and
Admin each read their own sub-directory and their own setting, so narrowing Admin does not
narrow Survey, and a language missing from Admin's sub-directory is not offered by the console
even when Survey offers it.

\begin{elicitnote}[Managing this from inside the applications is specified, not yet built]
Admin UC-027 specifies a \ui{Languages} screen for managing languages through the application
rather than the mounted directory and a startup property. It is not implemented in this release.
\end{elicitnote}

\subsection{Text direction and text size}
\label{sec:direction}

Arabic, Hebrew, Persian, Urdu, Pashto, Sindhi, Uyghur, Yiddish, Dhivehi and Kurdish (Sorani)
are laid out right-to-left by default. Every language is shown at the reader's own text size
unless something says otherwise. Three places can say otherwise, each winning over the one
before:
\begin{enumerate}
  \item The image's own manifest, \file{META-INF/i18n/i18n-config.json}.
  \item \file{i18n-config.json} at the root of the translations directory.
  \item A property for one tag on one application: \optlit{i18n.direction.<tag>} or
        \optlit{i18n.font-scale.<tag>} (Table~\ref{tab:common}).
\end{enumerate}
Both manifests declare \val{ar} right-to-left at a \opt{fontScale} of \val{1.15}:

\begin{elicitfile}
{ "locales": [ { "tag": "ar", "direction": "rtl", "fontScale": 1.15 } ] }
\end{elicitfile}

A site that mounts a language whose direction does not follow from its tag declares it in the
mounted \file{i18n-config.json}, or with a property when the change is for one application
only. \opt{direction} is \val{rtl} or \val{ltr}; no other direction is supported. The
applications' stylesheets mirror themselves for a right-to-left language, so no stylesheet per
language is needed.

\opt{fontScale} multiplies the reader's own base text size on pages shown in that language, and
changes nothing in any other. It exists because some scripts read smaller than Latin at the same
font size: Arabic does in the system fonts, so a page that is comfortable in English is cramped
in Arabic. The whole page scales in proportion (text, spacing and controls together), so
nothing is clipped, and a reader who has enlarged text in their browser keeps that enlargement
on top of the scale.

To raise or lower the scale, set \optlit{i18n.font-scale.ar} or edit the mounted manifest. The
change takes effect when each application that offers the language next restarts.

\begin{elicitnote}[Every entry is optional, and a bad value is ignored]
Each value is resolved on its own:
\begin{itemize}
  \item An entry or property may set the direction, the scale, or both.
  \item Anything a layer does not declare falls back to the layer below.
  \item A language no layer mentions renders exactly as it did before any of them existed.
  \item A malformed mounted manifest is logged and the image's own stands.
  \item A \opt{fontScale} that is not a number, or that falls outside \val{0.75} to \val{2.0},
        is logged and dropped, and that language is shown unscaled. It is never clamped into
        range: a size nobody asked for is worse than none.
\end{itemize}
\end{elicitnote}

\subsection{Which language a reader gets}

In order:
\begin{enumerate}
  \item a \opt{?lang=<tag>} parameter on the address, which an invitation link can carry;
  \item the language chosen earlier in the same browser session;
  \item the browser's preferred languages, matched against those offered;
  \item English.
\end{enumerate}
An unknown tag is ignored. The choice lasts for the session only; nothing is stored against a
respondent or a user.

\subsection{A language the curated directory does not have}
\label{sec:commission}

There are two routes:
\begin{itemize}
  \item Ask ElicitSoftware to add the language to the curated directory. It then arrives,
        reviewed, with a later release.
  \item Supply the language yourself, and have it from the next restart.
\end{itemize}
Either way the translation starts from the same document: each application's translation
request, \file{i18n/TRANSLATION_REQUEST.md} in its repository, which contains:
\begin{itemize}
  \item the application's purpose and audience;
  \item a glossary of terms that must stay consistent or untranslated;
  \item the placeholder and markup rules;
  \item every key with its English text, where it appears, its maximum length and any flags;
  \item the return format: one \file{translations_<tag>.properties}, with the same keys in the
        same order.
\end{itemize}

To supply a language yourself, do this for each of Survey and Admin:
\begin{enumerate}
  \item Take that module's \file{TRANSLATION_REQUEST.md} from the version you run.
  \item Hand it to a translator or an AI agent.
  \item Place the returned file in that module's sub-directory of your translations directory.
  \item Restart.
\end{enumerate}
Use the request from the version you run: one from another version lacks or misnames the texts
that changed, and those then appear in English. Nobody supporting your site can see a file you
supply or a curated file you edit, so keep it with the site's configuration, and send
corrections to ElicitSoftware so that the curated directory carries them too.

%% ================================================================== 12
\section{Custom reports}
\label{sec:reports}

A survey can end with reports: a summary of what the respondent said, a risk estimate, a
drawing. Elicit builds one of them itself, the \term{respondent report}, which echoes back the
answers a survey author marked for it. Every other report is a \term{report service}: a small
web application your site runs, which Survey and Admin call with a respondent's id and which
answers with the report.

This chapter is the contract between Elicit and a report service: when it is called, what it
is sent, what it must answer, and where it finds the answers to report on. It is written for
whoever builds the service, the \term{report developer}. Chapter~\ref{sec:fhhs} installs the
two report services that ship with Elicit, FHHS and Pedigree, which are the working example of
everything described here.

\subsection{How a report reaches a reader}

A survey carries an ordered list of report services. Each entry is a name, a description and
an address, and it arrives at a site in the survey's definition file. Two things ask for the
reports on that list:

\begin{description}
  \item[Survey,] when a respondent finishes. The respondent lands on a report page that calls
        every service on the list, in order, and shows each answer as a card headed with the
        report's name. A \ui{Generate PDF} button on the same page sets every report in type.
  \item[Admin,] when an administrator chooses \ui{Print reports} on a finished respondent's
        row. The console calls the same services and assembles their answers into one PDF.
\end{description}

\begin{elicitdiagram}{One report request. Survey and Admin send the same request and read the
same answer; only what they do with it differs.}
  \node[elicit app, minimum width = 2.7cm] (survey) at (0, 0.9)
    {Survey\\[-1pt]{\scriptsize report page}};
  \node[elicit app, minimum width = 2.7cm] (admin) at (0, -0.9)
    {Admin\\[-1pt]{\scriptsize \ui{Print reports}}};
  \node[elicit app, minimum width = 2.7cm, minimum height = 5.2em] (svc) at (6.4, 0)
    {Report service\\[-1pt]{\scriptsize your site's}};
  \node[elicit service, minimum width = 2.7cm, minimum height = 5.2em] (db) at (12.2, 0)
    {Reporting schema\\[-1pt]{\scriptsize\val{report_<survey>}}};
  \draw[elicit flow, {Stealth[length = 5pt]}-{Stealth[length = 5pt]}] (survey.east) --
    ([yshift = 0.55cm]svc.west);
  \draw[elicit flow, {Stealth[length = 5pt]}-{Stealth[length = 5pt]}] (admin.east) --
    ([yshift = -0.55cm]svc.west);
  \node[elicit label] at (3.7, 1.45) {request \val{{"id", "language"}}};
  \node[elicit label] at (3.7, -1.45) {answer \val{{"title", "innerHTML", "pdf"}}};
  \draw[elicit flow] (svc) -- node[elicit label, above = 2pt] {reads} (db);
\end{elicitdiagram}

A survey whose list is empty has no report page: a respondent who finishes goes straight to the
survey's post-survey address.

\subsection{The request}

Each service is sent an HTTP \val{POST} to exactly the address on the list, with a JSON body:

\begin{elicitfile}
{ "id": 1042, "language": "es-419" }
\end{elicitfile}

\begin{description}
  \item[\val{id}] The respondent's id: \val{survey.respondents.id}, which is also the
        \val{respondent_id} of every row the reporting schema holds about them.
  \item[\val{language}] The tag of the language the respondent answered in, so that the
        report can be written in it. It is \val{null}, or absent, when the respondent answered
        in the survey's base language, or the survey publishes no other.
\end{description}

\begin{plaintable}{tab:report-callers}{How each caller makes the request}{@{}>{\raggedright\arraybackslash}p{0.17\linewidth} >{\raggedright\arraybackslash}p{0.37\linewidth} >{\raggedright\arraybackslash}p{0.38\linewidth}@{}}{\thead{} & \thead{Survey} & \thead{Admin}}
When & The respondent reaches the report page after finishing, and on every later sign-in
with their access code. & An administrator prints the reports of a respondent whose status
is \textbf{Finished}. \\
\addlinespace
Language & The language the respondent is reading the survey in. & The language of the
respondent's most recent answer. \\
\addlinespace
Waits & 5 seconds to connect, 60 seconds for the answer. & The REST client's defaults. \\
\addlinespace
Credentials & None. & None. \\
\addlinespace
Uses & \val{innerHTML} for the page; \val{pdf} for \ui{Generate PDF}. & \val{pdf} only. \\
\end{plaintable}

\subsection{Where a report service runs}
\label{sec:reports-network}

A report service has no user interface. It serves no pages, has no sign-in and offers nothing
for a person to open in a browser; everything a reader sees of a report is what Survey and Admin
draw from its answer. A respondent's browser never calls a report service either: Survey calls
it on the respondent's behalf, and Admin on the administrator's.

So a report service is not exposed to the Internet. Run it on a private network where its
address answers only the computers running Survey and Admin: the network the containers share
under Compose, or a firewall rule that admits those hosts and no others. A service that another
report service calls, as FHHS calls Pedigree, answers only that service.

That isolation is why a report service needs no authentication. Neither caller presents a
credential, and the service checks none: the network has already decided who may ask. Build the
service on that assumption rather than adding a sign-in Survey and Admin cannot answer.

\begin{elicitsecurity}[The network is the only lock]
Anything that can reach a report service's address can ask it for any respondent's report. Never
publish its port beyond the Survey and Admin hosts, and never route it through the reverse proxy
that serves Survey to respondents. The shipped compose file publishes FHHS and Pedigree on the
evaluation host for local testing only; a real site removes those mappings.
\end{elicitsecurity}

\subsection{The answer}

A service answers \val{200} with a JSON object of three fields:

\begin{elicitfile}[fontsize = \footnotesize]
{
  "title": "Respondent Summary",
  "innerHTML": "<table class=\"card\"><tr><th>Age</th><td>47</td></tr></table>",
  "pdf": {
    "title": "Respondent Summary",
    "landscape": false,
    "pageBreak": false,
    "content": [
      { "text": "Answers given on 12 November 2026" },
      { "table": { "headers": ["Question", "Answer"], "widths": [3, 2],
                   "body": [["Age", "47"], ["Sex at birth", "Female"]] } },
      { "svg": "<svg xmlns=\"http://www.w3.org/2000/svg\" ...>...</svg>" }
    ]
  }
}
\end{elicitfile}

\begin{description}
  \item[\val{title}] The report's own name for itself. Survey heads the card with the
        report's name from the survey's list instead, translated into the respondent's
        language, so a site renames a report in the survey, not in the service.
  \item[\val{innerHTML}] An HTML fragment that Survey places, as it is, inside the report's
        card. A fragment, not a page: no \val{<html>} or \val{<body>}. An \val{<svg>} in it
        is centered in the card.
  \item[\val{pdf}] The same report as printable content, set in type by Admin's
        \ui{Print reports} and by Survey's \ui{Generate PDF}. It carries a \val{title}, whether
        the report is \val{landscape}, whether it starts on a new page (\val{pageBreak}), and
        a list of \val{content} blocks. Each block is one of a \val{text}, an \val{svg}
        drawing, or a \val{table} with \val{headers}, relative column \val{widths} and a
        \val{body} of rows. A block carrying a drawing is set on a landscape page.
\end{description}

\begin{elicitcaution}[Without \val{pdf} the console prints nothing]
Admin's printout is made from \val{pdf} alone. A service that answers only \val{innerHTML}
shows on the respondent's report page, but is left out of the administrator's PDF without a
word.
\end{elicitcaution}

\begin{elicitcaution}[The HTML is trusted]
Survey does not sanitize \val{innerHTML}; it is part of Survey's own page. Build it from the
answers with every value escaped, and only ever list services your site controls.
\end{elicitcaution}

When a service fails, answers with an error or does not answer in time, Survey shows an error
card in that report's place and Admin prints an error block for it. The other reports are
shown and printed as usual, and the respondent's survey is finished either way.

\subsection{Where the answers are}
\label{sec:reports-data}

A report service reads the survey's \term{reporting schema}, not the transactional tables.
Survey builds one per installed survey when the definition is applied
(Chapter~\ref{sec:definition}). It holds the answers to the questions the survey's author
tagged for analysis, one row per section a respondent reached, with each analysis tag as a
column.

\subsubsection*{Find the schema on every request}

The schema's name belongs to the site: it defaults to \val{report_} followed by the survey's
name, and an administrator can rename it at any time. Look it up by the survey's key, which
is the same at every site, on every request, rather than writing the name into the service:

\begin{elicitsql}
SELECT report_schema FROM survey.surveys WHERE survey_key = :survey_key;
\end{elicitsql}

A service that does this follows a rename on its next request. Treat the result as an
identifier, quoted as one, never as text spliced into a query.

\subsubsection*{Read one respondent}

\begin{plaintable}{tab:report-views}{The objects a report reads, in the survey's reporting schema}{@{}>{\raggedright\arraybackslash}p{0.27\linewidth} >{\raggedright\arraybackslash}p{0.67\linewidth}@{}}{\thead{Object} & \thead{What a row is}}
\val{fact_sections_view} & One section a respondent reached; a repeated section has a row per
repetition. Filter on \val{respondent_id}. Fixed columns \val{step}, \val{step_instance},
\val{section}, \val{section_instance}, \val{question} and \val{item}, then one column per
analysis tag holding that respondent's answer (for the Family Health History Survey,
\val{relationship}, \val{age}, \val{gender} and the rest). \\
\addlinespace
\val{fact_respondents_view} & One respondent: \val{id} (the respondent's id), \val{status},
and the dates they were registered, first signed in and finished. \\
\addlinespace
\val{dim_step}, \val{dim_section}, \val{dim_<tag>} & The values each column can hold, for a
report that lists the options a respondent did \emph{not} choose. \\
\end{plaintable}

\begin{elicitsql}
SELECT step, section, relationship, age, gender
  FROM report_family_history_survey.fact_sections_view
 WHERE respondent_id = :id
 ORDER BY step, step_instance, section_instance;
\end{elicitsql}

The fact rows of a respondent are written when they finish, so a report is complete only for a
finished respondent, which is exactly when both callers ask for one. The schema holds answers,
not identity: a subject's name and contact details are not in it.

\subsubsection*{Connect as the reporting role}

The role script creates \val{surveyreport_user} (Chapter~\ref{sec:database}) for exactly this.
Survey grants it read access to every reporting schema it builds, and, so that it can find
them, to the two columns of \val{survey.surveys} that name each survey's schema. That grant is
this statement, which Survey's migrations run for you:

\begin{elicitsql}
GRANT SELECT (survey_key, report_schema) ON survey.surveys
   TO surveyreport_user;
\end{elicitsql}

Run it yourself, as the schema owner, if your report service connects as a role of its own
rather than \val{surveyreport_user}. It is granted column by column, so the role still reads
nothing else of the transactional schema.

FHHS connects as \val{survey_user} instead, because it also reads the subject's details for the
document it files elsewhere (Chapter~\ref{sec:fhhs}). A report that only reports needs no more
than the reporting role.

\subsection{A minimal report service}

The contract is small enough to meet in any language. This one, in Python with Flask and
psycopg, lists a respondent's answers to two tags:

\begin{elicitcmd}[fontsize = \footnotesize]
import os, html
import psycopg
from psycopg import sql
from flask import Flask, request, jsonify

app = Flask(__name__)
SURVEY_KEY = os.environ["SURVEY_KEY"]  # the survey's key, the same at every site
DB = os.environ["REPORT_DB"]  # postgresql://surveyreport_user:...@db/survey

@app.post("/report")
def report():
    respondent = int(request.get_json()["id"])
    with psycopg.connect(DB) as conn:
        schema = conn.execute(
            "SELECT report_schema FROM survey.surveys WHERE survey_key = %s",
            (SURVEY_KEY,)).fetchone()[0]
        rows = conn.execute(
            sql.SQL("SELECT relationship, age FROM {}.fact_sections_view"
                    " WHERE respondent_id = %s ORDER BY step_instance")
               .format(sql.Identifier(schema)), (respondent,)).fetchall()
    body = [[str(r or ""), str(a or "")] for r, a in rows]
    cells = "".join(f"<tr><td>{html.escape(r)}</td><td>{html.escape(a)}</td></tr>"
                    for r, a in body)
    return jsonify(
        title="Relatives",
        innerHTML=f"<table><tr><th>Relative</th><th>Age</th></tr>{cells}</table>",
        pdf={"title": "Relatives",
             "content": [{"table": {"headers": ["Relative", "Age"],
                                    "widths": [1, 1], "body": body}}]})
\end{elicitcmd}

Run it beside the applications, on the same network, and give it an address Survey and Admin
can reach and nothing outside that network can (Section~\ref{sec:reports-network}), such as
\val{http://relatives:8080/report}.

\subsection{Adding the service to a survey}

A report service is registered in the survey, never in a site's database:

\begin{enumerate}
  \item In Author, open the survey's \ui{Reports\ldots} page and choose
        \ui{Add report service\ldots}.
  \item Give it a \ui{Name}, which heads the report card, a \ui{Description}, and the
        \ui{Address} the applications will call. Translate the name and description if the
        survey publishes other languages.
  \item Drag it to its place in the list, and export the survey.
  \item Apply the exported file at the site through Admin's \ui{Apply Survey Definition}
        (Chapter~\ref{sec:definition}).
  \item Open \ui{System} $>$ \ui{Connections} in Admin. Each report service is listed by
        name and survey; check it to confirm it answers.
\end{enumerate}

\begin{elicitnote}[The list in the file is the list at the site]
Every apply replaces the survey's report list with the one in the file, and removes any service
the file no longer names. An address changed by hand in the database therefore lasts only until
the next apply. Change it in Author and export again.
\end{elicitnote}

The \ui{Connections} check sends each service one read-only request and reports whether
anything answered; it never asks for a report. A service that answers that request with
\val{405}, as one accepting only \val{POST} does, is reachable.

\subsection{Doing something when a respondent finishes}
\label{sec:psa}

A report is drawn when somebody asks for it. Some sites also need something to happen at the
moment a respondent finishes: a document filed in a records system, a message sent to a
study coordinator. That is a \term{post-survey action}.

A survey carries an ordered list of them beside its reports, and like the reports they arrive
in the definition file. An author lists them on the survey's \ui{Post-survey actions\ldots} page
in Author: each a name, a description, a method, an address and whether it sends the subject's
External ID. Like the report list, the file's list is the site's list: every apply inserts,
updates and deletes actions to match it, and an action it deletes takes its delivery records with
it.

\subsubsection*{The call}

When a respondent finishes, Survey calls each action in order, waiting 5 seconds to connect and
60 for the answer:

\begin{description}
  \item[Method] \val{POST}, \val{GET} or \val{PUT}, as the author chose.
  \item[Address] As written, after its placeholders are filled and URL-encoded:
        \texttt{\small\{respondentId\}} with the respondent's id, and, on an action that sends the
        External ID, \texttt{\small\{externalId\}} with the subject's External ID.
  \item[Body] For \val{POST} and \val{PUT}, \texttt{\small\{"id": <respondent id>\}}. An action
        that sends the External ID adds it, and the department's code, because an External ID is
        unique only within its department:
        \texttt{\small\{"id": 1042, "externalId": "MRN00123", "departmentCode": "CARD"\}}. A
        \val{GET} has no body.
  \item[Answer] Any \val{2xx} is success. The body of the answer is never read: unlike a report
        service, an action returns no HTML and no \val{pdf}, and owes Survey no particular
        response.
\end{description}

\begin{elicitsecurity}[The External ID travels only over https]
An External ID usually identifies a person in another system, such as a medical record number.
An action that sends it must use an \val{https} address. Author will not save one at
\val{http}, Admin refuses a definition file carrying one, and Survey refuses to make the call,
so even a hand-edited file cannot send it in the clear.
\end{elicitsecurity}

\subsubsection*{Giving an action a credential}

Unlike a report service, a post-survey action may call a service outside the site, which may
want a credential. The credential belongs to the site, so it is set in Survey's configuration,
keyed by the action's key, and never travels in a definition file:

\begin{plaintable}{tab:psa-credentials}{Credentials for one post-survey action, on Survey}{@{}>{\raggedright\arraybackslash}p{0.47\linewidth} >{\raggedright\arraybackslash}p{0.47\linewidth}@{}}{\thead{Setting} & \thead{What it sends}}
\optlit{elicit.post-survey-action."<key>".bearer-token} & \texttt{\small Authorization: Bearer <token>}, the token as given. Supply it as an environment variable or a secret, never in a file under version control. \\
\end{plaintable}

\val{<key>} is the action's key, the second field of its record in the definition file and the
same at every site. The token is one the service issued to your site; Survey presents it as
given and never requests or refreshes one, so when the service retires a token, set the new one
and restart Survey. An action without a token is called without a credential. Survey's startup
diagnostics say whether each action has a credential, never its value.

\subsubsection*{When a call fails}

Survey records each delivery in \val{survey.respondent_psa}, and retries a call that failed:

\begin{plaintable}{tab:psa-states}{Delivery states Survey records}{@{}>{\raggedright\arraybackslash}p{0.20\linewidth} >{\raggedright\arraybackslash}p{0.74\linewidth}@{}}{\thead{State} & \thead{Meaning}}
\val{PENDING}, \val{RESENDING} & Being called, for the first time or again. \\
\val{DELIVERED} & The service answered \val{2xx}. \\
\val{CALL_FAILED} & No answer in time, no connection, or another status; retried. \\
\val{CALL_REFUSED} & Refused before the call: an External-ID action at an \val{http} address, an
address that needs an External ID the subject does not have, or a placeholder it may not use.
Never retried, because a retry cannot change the outcome. \\
\end{plaintable}

Every \opt{elicit.post-survey-action.retry.interval} (15 minutes unless changed), Survey calls
each \val{CALL_FAILED} delivery again, until it is delivered or has had
\opt{elicit.post-survey-action.retry.max-attempts} attempts (24, the first included). Each
retried delivery is claimed first, so several Survey instances never call it twice, and
\opt{elicit.post-survey-action.retry.enabled} turns the schedule off (Table~\ref{tab:survey}).

\begin{elicitcaution}[Survey retries the call, not the service's work]
Once a call is answered, what the service does next is its own. If it records its own progress
on the same delivery, as FHHS does with \val{STARTED}, \val{COMPLETED} and \val{FAILED}, Survey
never overwrites that and never retries it. A service that must finish reliably answers at once,
does the work in the background, and keeps its own record and its own retries, as FHHS does for
the documents it files (Section~\ref{sec:fhhs-sftp}).
\end{elicitcaution}

%% ================================================================== 13
\section{The Family History reports}
\label{sec:fhhs}

The Family Health History Survey ends with three reports, all answered by one report service,
\ui{FHHS}, which asks a second, \ui{Pedigree}, to draw the family tree. Together they are the
reference implementation of the contract in Chapter~\ref{sec:reports}: a site that runs the
survey installs them from this chapter, and a report developer reads them as a working
example.

Both are optional. A site that does not run the Family Health History Survey runs neither.

\begin{elicitdiagram}{The Family History reports. FHHS answers three report services and one
post-survey action; Pedigree only ever hears from FHHS.}
  \node[elicit app, minimum width = 2.4cm] (survey) at (0, 1.0) {Survey};
  \node[elicit app, minimum width = 2.4cm] (admin) at (0, -1.0) {Admin};
  \node[elicit app, align = left, minimum width = 4.2cm, inner ysep = 7pt] (fhhs) at (5.6, 0)
    {\textbf{FHHS}\\[3pt]
     \texttt{\scriptsize /proband/report}\\[-1pt]
     \texttt{\scriptsize /casummary/report}\\[-1pt]
     \texttt{\scriptsize /pedigree/report}\\[-1pt]
     \texttt{\scriptsize /familyhistory/generate}};
  \node[elicit app, minimum width = 2.6cm] (ped) at (11.4, 1.0)
    {Pedigree\\[-1pt]\texttt{\scriptsize /svg}};
  \node[elicit service, minimum width = 2.6cm] (sftp) at (11.4, -1.0)
    {SFTP server\\[-1pt]{\scriptsize optional}};
  \node[elicit service, minimum width = 4.2cm] (db) at (5.6, -2.7)
    {\texttt{\footnotesize report\_family\_history\_survey}};
  \draw[elicit flow] (survey.east) -- node[elicit label, above = 2pt, sloped] {reports, action}
    ([yshift = 0.5cm]fhhs.west);
  \draw[elicit flow] (admin.east) -- node[elicit label, below = 2pt, sloped] {reports}
    ([yshift = -0.5cm]fhhs.west);
  \draw[elicit flow] ([yshift = 0.5cm]fhhs.east) -- node[elicit label, above = 2pt, sloped]
    {pedigree table} (ped.west);
  \draw[elicit flow] ([yshift = -0.5cm]fhhs.east) -- node[elicit label, below = 2pt, sloped]
    {PDF + index} (sftp.west);
  \draw[elicit flow] (fhhs) -- node[elicit label, right = 2pt] {reads} (db);
\end{elicitdiagram}

\subsection{What they answer}

\begin{plaintable}{tab:fhhs-endpoints}{FHHS's addresses, as the survey's definition lists them}{@{}>{\raggedright\arraybackslash}p{0.30\linewidth} >{\raggedright\arraybackslash}p{0.17\linewidth} >{\raggedright\arraybackslash}p{0.45\linewidth}@{}}{\thead{Address} & \thead{Listed as} & \thead{What it answers}}
\val{/proband/report} & Report \emph{Respondent} & The respondent's own history, as a table. \\
\addlinespace
\val{/casummary/report} & Report \emph{Cancer Summary} & Each relative's cancers and ages at
diagnosis. \\
\addlinespace
\val{/pedigree/report} & Report \emph{Family Pedigree} & The family tree, drawn by Pedigree,
as an SVG on a landscape page of its own. \\
\addlinespace
\val{/familyhistory/generate} & Post-survey action & Files all three as one PDF on an SFTP
server, when the site turns that on (Section~\ref{sec:fhhs-sftp}). \\
\end{plaintable}

Each report reads the survey's \val{fact_sections_view} for the respondent, resolving the
schema's name from \val{survey.surveys.report_schema} on every request, exactly as
Section~\ref{sec:reports-data} describes. The source is in the FHHS repository, one package
per report (\file{proband}, \file{casummary}, \file{pedigree}); it is the place to start a
report of your own.

\subsection{Installing with Docker Compose}

The shipped compose file already runs both, as the services \val{fhhs} and \val{pedigree}.
FHHS waits for Admin to be ready before it starts (Chapter~\ref{sec:topology} explains why), and
connects with the same database passwords as Survey (Chapter~\ref{sec:secrets}). Bring the
stack up as in Chapter~\ref{sec:compose}, then:

\begin{enumerate}
  \item Watch FHHS report \emph{not ready}. Until the survey is installed it logs one warning
        naming the survey it is waiting for, its readiness check gives the same reason, and
        it refuses report requests with \val{503} and that message.
  \item Make sure the survey's report addresses reach FHHS from the Survey and Admin containers
        (see \emph{The report addresses} below).
  \item Apply \file{FHHS/family-history-survey.elicit} in Admin
        (Chapter~\ref{sec:definition}).
  \item FHHS becomes ready on its next readiness probe. Nothing needs restarting.
  \item In Admin, open \ui{System} $>$ \ui{Connections} and check the three report services and
        the post-survey action.
  \item Walk the survey to the end as a respondent. The report page shows three cards, the
        last of them the pedigree.
\end{enumerate}

\subsection{Installing without Compose}

\begin{enumerate}
  \item Start \file{elicitsoftware/pedigree} at the release's tag. It listens on \val{8080},
        reads no configuration and keeps no data. Its health endpoint is \opt{/healthz}.
  \item Start \file{elicitsoftware/fhhs} at the same tag, \emph{after Admin is ready}, with the
        datasource settings Survey uses and the settings in Table~\ref{tab:fhhs}. At minimum:
\end{enumerate}

\begin{elicitcmd}[fontsize = \footnotesize]
docker run -d --name fhhs \
  -e QUARKUS_DATASOURCE_JDBC_URL='jdbc:postgresql://db.example.org:5432/survey' \
  -e QUARKUS_DATASOURCE_OWNER_JDBC_URL='jdbc:postgresql://db.example.org:5432/survey' \
  -e ELICIT_DB_PASSWORD='<survey_user password>' \
  -e ELICIT_OWNER_DB_PASSWORD='<elicit_owner password>' \
  -e PEDIGREE_URL='http://pedigree.internal:8080/svg' \
  elicitsoftware/fhhs:<version>
\end{elicitcmd}

Keep the connection options of the shipped JDBC address (Appendix~\ref{sec:config}) when you
point FHHS at your own server. FHHS reports ready at \opt{/q/health/ready} once the survey is
applied, as above.

\subsection{The report addresses}

The survey's definition file lists FHHS's three reports and its action at
\val{http://host.docker.internal:8082/…}: FHHS's published port on the evaluation host, which
the Survey and Admin containers reach through Docker Desktop's host alias. That works for the
shipped compose file on a developer's machine and nowhere else.

A real site gives the four addresses the name the applications know FHHS by on their own
network, for example \val{http://fhhs:8080/proband/report}. Change them in Author, on the
survey's \ui{Reports\ldots} page for the three reports and on its \ui{Post-survey actions\ldots}
page for the action, and export the survey, then apply that file instead of the shipped one.
A change made in the database is undone by the next apply.

\subsection{FHHS settings}

FHHS takes the datasource and observability settings Survey does, with the same names and
passwords (Appendix~\ref{sec:config}). Its own are these.

\begin{settings}{tab:fhhs}{FHHS settings}
\opt{pedigree.url} & --- & \chgset & Pedigree's drawing endpoint, for example
\val{http://pedigree:8080/svg}. No value is shipped; the compose file sets it.
Environment: \opt{PEDIGREE_URL}. \\
\opt{family.history.survey.key} & \val{5e91c606-59a1-450a-a8d7-2f1530ff472b} & \chgkeep & The
key of the survey FHHS serves. It is the key in the shipped definition file, the same at every
site. \\
\opt{family.history.upload.psa.key} & \val{9f57b644-fecd-557d-a466-07ec19bb95c4} & \chgkeep &
The key of the post-survey action FHHS records its uploads against. \\
\opt{family.history.upload.psa.id} & --- & \chgmay & The same action by its id at this site.
Set, it overrides the key; leave it unset. \\
\opt{family.history.async.threads} & \val{2} & \chgmay & How many family-history documents
FHHS builds at once. \\
\opt{app.base.url} & \val{http://localhost:8080} & \chgmay & The survey's public address,
printed at the foot of the filed PDF. Set it when filing is on. Environment:
\opt{APP_BASE_URL}. \\
\opt{family.history.sftp.enabled} & \val{false} & \chgmay & Whether the post-survey action
files documents on an SFTP server (Section~\ref{sec:fhhs-sftp}). Left \val{false}, the
action answers and files nothing. \\
\opt{family.history.sftp.host} & --- & \chgset & The SFTP server, when filing is on. \\
\opt{family.history.sftp.port} & \val{22} & \chgmay & \\
\opt{family.history.sftp.path} & --- & \chgset & The directory documents are written to. \\
\opt{family.history.sftp.username} & --- & \chgset & The account FHHS signs in as. \\
\opt{family.history.sftp.password} & --- & \chgmay & Its password, when the server accepts
one. A secret. \\
\opt{family.history.sftp.privateKey} & --- & \chgmay & A private key instead: a path to a
mounted key file, or the key itself. Takes precedence over the password. A secret. \\
\opt{family.history.sftp.knownHosts} & --- & \chgset & The server's host key, as
\val{ssh-keyscan} prints it. FHHS refuses a server whose key it cannot verify. \\
\opt{family.history.sftp.strictHostKeyChecking} & \val{true} & \chgkeep & Whether to insist on
that verification. The compose file turns it off for the bundled server, which makes a new key
on every start; never turn it off for a real one. \\
\opt{family.history.sftp.timeout} & \val{30000} & \chgmay & How long, in milliseconds, to wait
for the server. \\
\opt{family.history.sftp.xml.template} & an index document & \chgmay & The index written
beside each PDF (Section~\ref{sec:fhhs-sftp}). \\
\end{settings}

\subsection{Filing the family history elsewhere}
\label{sec:fhhs-sftp}

When \opt{family.history.sftp.enabled} is \val{true}, FHHS answers the survey's post-survey
action by building one PDF of every report on the survey's list and uploading it to the SFTP
server, beside an index document a records system can file it by:

\begin{elicitfile}
<path>/<xid>.pdf
<path>/<xid>-index.xml
\end{elicitfile}

The index comes from \opt{family.history.sftp.xml.template}, in which these placeholders are
replaced from the respondent's record: \val{{RespondentId}}, \val{{Xid}},
\val{{SurveyId}}, \val{{FirstName}}, \val{{MiddleName}}, \val{{LastName}},
\val{{Dob}}, \val{{Email}}, \val{{Phone}}, \val{{DepartmentName}},
\val{{DepartmentID}}, \val{{AccessCode}}, \val{{Status}}, \val{{Created}} and
\val{{Finalized}}. The compose file carries a commented example.

FHHS answers the action at once and does the work in the background. It records every upload
against the respondent in \val{survey.respondent_psa}, as started, completed or failed, and
every fifteen minutes retries the failed ones, up to fifty attempts. That is the pattern
Section~\ref{sec:psa} recommends for any action that must arrive. The two retries never
overlap: Survey calls FHHS again only when its call never reached FHHS (\val{CALL_FAILED});
once FHHS has answered, the filing and its retries are FHHS's.

\subsection{Pedigree}

Pedigree is an R service built on the \emph{Pedixplorer} package. The endpoint that matters is
\val{POST /svg}. It takes a tab-separated pedigree table, one row per family member, as the
multipart field \val{ped}, and answers with the drawing as an SVG. The table's columns are
\val{famid}, \val{id}, \val{sex}, \val{dadid} and \val{momid}, and optionally \val{deceased},
\val{proband}, \val{affection}, \val{avail}, \val{display_id} and \val{cancer_label}. FHHS
builds the table from the respondent's answers and waits 15 seconds for the drawing.
\val{GET /healthz} answers \val{{"status":"ok"}} when the service is up.

Pedigree is a report service's helper, not a report service: nothing in Elicit calls it but
FHHS, so it never appears on a survey's report list.

\elicitpart{Running the site}{Checking that every part of the installation works, what to do
when one does not, and where to read further.}

%% ================================================================== 14
\section{Verifying the installation}
\label{sec:verify}

Work through these checks in order: each one relies on the one before it.

\begin{enumerate}
  \item \textbf{Every service is ready.} \file{./status.sh} lists the containers and asks each
        application whether it is ready: \opt{/q/health/ready} for Survey, Admin, FHHS and
        Author, \opt{/healthz} for Pedigree. A service that is up but not ready is reported as
        such, with the reason its readiness check gives. Under your own container platform,
        ask each \opt{/q/health/ready} directly. For Survey, ready means its migrations have
        finished.
  \item \textbf{The schema exists.} The \val{survey} schema carries tables, and each module's
        migration history is present.
  \item \textbf{The reporting schema is built.} After a definition is applied, the survey's
        \val{report_schema} is set and \val{dim_step} in that schema has rows
        (Chapter~\ref{sec:definition}).
  \item \textbf{Sign-in works.} Sign in to Admin. A redirect that fails here is almost always
        the identity provider's issuer address (Chapter~\ref{sec:oidc}).
  \item \textbf{The console agrees.} Admin's \ui{System} section reports on the database,
        branding, mail, the reachability of every outbound dependency (report services
        included) and the current sign-in's roles. Its screens are shown below.
  \item \textbf{A respondent can answer.} Register a subject, issue an access code, and walk
        the survey to the end. If the survey has reports, the report page shows each of them.
\end{enumerate}

\screenshot{11-system-overview}{The System overview page. It lists the setup work still
outstanding, which on a finished installation makes it the shortest possible verification.}

\screenshot{12-system-database}{The System database page: both connections, and whether each
module's migrations are in place.}

\screenshot{14-system-email}{The System email page: the effective settings, and a button that
sends a real message through them.}

\screenshot{15-system-connections}{The System connections page: a bounded reachability check of
every outbound dependency, each of the survey's report services among them.}

\screenshot{16-system-oidc}{The System OIDC page: the current sign-in's roles, and its tokens
masked. Secrets are reported as present or absent, never shown.}

\begin{elicitnote}[Survey writes a diagnostics report to its log]
A few seconds after startup Survey writes a startup diagnostics report, which also checks
every report service it will call. Anything an operator must act on is logged as a warning;
the rest is logged at information level, and its category is pinned so the report appears even
when the overall log level is \val{WARN}. It is the first place to look after a start that
seemed to work.
\end{elicitnote}

%% ================================================================== 15
\section{When it does not work}
\label{sec:trouble}

\begin{plaintable}{tab:failures}{Failures that are actually seen}{@{}>{\raggedright\arraybackslash}p{0.29\linewidth} >{\raggedright\arraybackslash}p{0.31\linewidth} >{\raggedright\arraybackslash}p{0.33\linewidth}@{}}{\thead{Symptom} & \thead{Cause} & \thead{Fix}}
A module fails its migration on a missing table or sequence. & It started before Survey
migrated a new database. & Let Survey become ready, then start the module again. Nothing is
corrupted. \\
\addlinespace
FHHS is not ready and answers report requests with \val{503}, naming a survey. & Its survey
has not been imported yet, or its reporting schema has not been built. & Apply the definition,
or apply it again to build the schema (Chapter~\ref{sec:definition}). FHHS goes healthy on its
next probe; nothing needs restarting. \\
\addlinespace
A report card shows an error, or a report service is unreachable on \ui{Connections}. & The
address in the survey's report list does not resolve from inside the Survey and Admin
containers, or the service is down. & Correct the address in Author and apply the new export
(Chapter~\ref{sec:reports}); a fix in the database lasts only until the next apply. \\
\addlinespace
A report shows on the respondent's page but is missing from the console's printout. & The
service answers no \val{pdf}. & Return the printable description as well
(Chapter~\ref{sec:reports}). \\
\addlinespace
A post-survey action's deliveries stay \val{CALL_REFUSED}. & The action sends the External ID
to an \val{http} address, uses a placeholder it may not, or needs an External ID the subject
lacks; the delivery's \val{error_msg} says which. & Correct the action in Author and apply the new
export, or give the subject an External ID. A refused delivery is not retried
(Section~\ref{sec:psa}). \\
\addlinespace
A post-survey action's calls fail with \val{401} or \val{403}. & No token reached the
service, or the service no longer accepts it: the setting's name does not hold the action's key,
or the token has expired or been revoked. & Check the line for the action in Survey's startup
diagnostics, which says whether it has a credential; if it does, ask the service for a new token
(Section~\ref{sec:psa}). \\
\addlinespace
Sign-in redirects fail, or land on the wrong service. & The issuer address does not resolve
to the same provider from the browser and from inside the container. & Use one address that
works from both (Chapter~\ref{sec:oidc}). \\
\addlinespace
A service is reported unhealthy although it responds in a browser. & The health check was
pointed at the published host port. & Point it at \val{8080} inside the container. \\
\addlinespace
The database's first start after a reset fails on wrong ownership of its data directory. & A
known file-sharing fault of Docker Desktop for Mac (\val{docker/for-mac\#7415}), not a broken
copy. & Run \file{resetDatabase.sh} again and restart. If it recurs, switch Docker Desktop's
file sharing to gRPC FUSE. \\
\addlinespace
The console shows a banner about default accounts on every screen. & The seeded accounts
\val{admin} and \val{user} have not been renamed. & Rename both
(Chapter~\ref{sec:firstsignin}). \\
\addlinespace
A dialog blocks the console and will not close. & No department exists. & Follow its link and
create one (Chapter~\ref{sec:firstsignin}). \\
\addlinespace
A signed-in user can only log out. & They belong to no department. & Assign them one in
\ui{Users}. \\
\addlinespace
A mounted language is never offered. & Any one of these:
\begin{itemize}
  \item The directory is not mounted on that application, or \opt{i18n.file.system.path} does
        not name the mount point.
  \item The file is in the other application's sub-directory, or its name uses a hyphen
        (\file{translations_es-419.properties}).
  \item The tag is missing from \opt{i18n.offered.locales}.
  \item The module has not been restarted since the change.
\end{itemize}
& The application logs a directory it found nothing in, and each listed tag with no file. Fix
the cause (Section~\ref{sec:language}) and restart. \\
\end{plaintable}

\subsection{Upgrading an existing site}

Upgrade procedures belong to a release, and this manual, stamped with one version, may outlive
it. They are kept in \file{DeploymentScript.md} in the umbrella repository, which covers,
among others, the move to Kimball Type 2 slowly changing dimensions and the rename of the
respondent credential to the access code. Follow the procedure for the release you are moving
to, not this manual.

\begin{elicitcaution}[Upgrade the modules that share a schema together]
Survey, Admin and FHHS share the \val{survey} schema and migrate on top of each other.
Upgrading one without the others is what breaks. In the compose file, \opt{depends_on} orders
\emph{startup} only; it does not check again a dependency that is already running.
\end{elicitcaution}

%% ================================================================== 16
\section{Where to look next}

\begin{description}
  \item[The administrator's manual] Running the console: departments, users, subjects, access
        codes, invitations, reports, and applying survey definitions.
  \item[The Author installation manual] Provisioning the authoring tool: its database, its
        identity-provider client and roles, and the preview Survey it runs beside.
  \item[\file{DeploymentScript.md}] Release-specific upgrade procedures, and the post-install
        template data for the Family Health History Survey.
  \item[The FHHS repository] The source of the three Family History reports, one package per
        report, and its use cases under \file{docs/use_cases/}: the most complete example of
        a report service there is.
  \item[\file{docs/BRAND_SYSTEM_IMPLEMENTATION_GUIDE.md}] How the brand system is built, for
        someone extending it rather than mounting one.
  \item[\file{docs/I18N_IMPLEMENTATION_GUIDE.md}] The same, for translations.
  \item[\url{https://quarkus.io/guides/all-config}] Every configuration setting this manual
        does not list.
\end{description}

\vspace{1em}

\begin{elicitnote}[Which version is this?]
This manual describes the Elicit platform \textbf{\elicitVersion}, built
\textbf{\elicitBuildDate}. The same pair is printed at the foot of every page, so a printed
copy can always be matched back to the release it came from.
\end{elicitnote}

\elicitpart{Reference}{Every setting Survey and Admin define, every Quarkus setting their images
set, and how a setting reaches an application.}
\appendix
\renewcommand{\thesection}{\Alph{section}}

%% ================================================================== A
\section{Configuration reference}
\label{sec:config}

This appendix lists every setting Survey and Admin define and every Quarkus setting their
images set. FHHS's own settings are with its installation, in Table~\ref{tab:fhhs}; it shares
the datasource and observability settings below.

\subsection{How configuration reaches an application}

Every module is a Quarkus application, so every setting can be given as a property or as an
environment variable, whichever your deployment makes easier. When a setting is given in more
than one place, the source with the highest ordinal wins:

\begin{plaintable}{tab:sources}{Where a setting can come from}{@{}>{\raggedright\arraybackslash}p{0.62\linewidth} >{\raggedleft\arraybackslash}p{0.12\linewidth}@{}}{\thead{Source} & \thead{Ordinal}}
System properties, passed on the command line & 400 \\
Environment variables & 300 \\
Environment variables from a \file{.env} file in the working directory & 295 \\
\file{application.properties} in a \file{config} directory, relative to the application's
working directory in its container & 260 \\
The \file{application.properties} inside the image: the defaults this appendix lists & 250 \\
\end{plaintable}

A \file{.env} file therefore sits \emph{below} real environment variables, not above them.

\subsubsection*{From a property name to an environment variable}

Upper-case the name and replace every character that is neither a letter, a digit nor an
underscore with an underscore. So \opt{quarkus.datasource.jdbc.url} becomes
\opt{QUARKUS_DATASOURCE_JDBC_URL}, and a quoted segment produces a doubled underscore on each
side: \opt{quarkus.log.category."io.agroal".level} becomes
\opt{QUARKUS_LOG_CATEGORY__IO_AGROAL__LEVEL}.

\begin{elicitcaution}[The mapping is lossy]
Names that differ only in punctuation collide: \opt{foo.bar-baz} and \opt{foo.bar.baz} both
become \opt{FOO_BAR_BAZ}. Where that matters, give the setting as a property instead. In a
container you can also give the dotted name directly as an environment entry, which avoids
the question entirely and is what the shipped compose file does.
\end{elicitcaution}

\subsubsection*{Profiles}

A name prefixed with a profile, such as \optlit{\%dev.quarkus.http.port}, applies only under that
profile. Three prefixes appear in Survey's and Admin's files, and none of them configures a
deployment:

\begin{itemize}
  \item \optlit{\%dev.}: values for \val{quarkus:dev} on a developer's machine.
  \item \optlit{\%test.}: values for the module's own test run.
  \item \optlit{\%docker.}: a build-time profile that does one thing, switch the container-image
        build on. It sets nothing a running deployment reads.
\end{itemize}

There are no \optlit{\%prod.} entries. Everything a deployment runs with is either an unprefixed
default from the tables below or a value your site supplies.

\subsection{Reading the tables}

The \textbf{Change} column says what the setting is to you:

\begin{description}
  \item[\chgset] Your site must set it. The shipped value is a development value or a
        placeholder and is wrong anywhere else.
  \item[\chgmay] Your site may set it. The shipped value is a working default.
  \item[\chgkeep] Leave it alone. It is an internal default that happens to live in the same
        file, and changing it changes how the application works.
\end{description}

\subsection{Settings common to Survey and Admin}

\begin{settings}{tab:common}{Settings Survey and Admin both read}
\opt{brand.file.system.path} & \val{/brand} & \chgmay & The mounted brand directory. The
compose file sets \val{/opt/brand}. See Chapter~\ref{sec:brand}. \\
\opt{brand.local.path} & \val{brand} & \chgkeep & A brand directory beside the running
application, searched between the mount and the embedded default. \\
\opt{i18n.file.system.path} & \val{/i18n} & \chgmay & The mounted translations directory. Each
application reads its own sub-directory of it (\opt{i18n.app.name}). The compose file mounts
\file{elicit-i18n/} at \val{/opt/i18n} and sets this to match. A directory that does not exist is
logged and the application offers English only. See Chapter~\ref{sec:i18n}. \\
\opt{i18n.app.name} & per app & \chgkeep & This application's sub-directory of the translations
directory: \val{survey} or \val{admin}. Must be a plain directory name. \\
\opt{i18n.offered.locales} & --- & \chgmay & The languages this application offers, when it
should offer fewer than are mounted: a comma-separated list of tags. Unset offers every mounted
language. English is always offered. A listed tag with no mounted file is logged and not
offered, so this cannot add a language. A language left out is hidden from the selector, a
\opt{?lang=} naming it is refused, and survey content is not served in it either, because
content is only offered in a language the site also offers for the chrome. See
Chapter~\ref{sec:i18n}. \\
\optlit{i18n.direction.<tag>} & --- & \chgmay & Overrides both manifests (the image's
\file{META-INF/i18n/i18n-config.json} and the \file{i18n-config.json} at the root of the
translations directory) for one language tag: \val{rtl} or \val{ltr}. Both manifests declare
\val{ar} right-to-left; every other language is inferred from its tag. Any other value is logged
and ignored. See Section~\ref{sec:direction}. \\
\optlit{i18n.font-scale.<tag>} & --- & \chgmay & Overrides both manifests' \val{fontScale} for
one tag, multiplying the root font size on pages shown in that language. Both manifests scale
\val{ar} to \val{1.15}; every other language is \val{1.0}. Must be a number from \val{0.75} to
\val{2.0}; a value outside that range is logged and dropped rather than clamped. See
Section~\ref{sec:direction}. \\
\opt{i18n.pseudo-locale.enabled} & \val{false} & \chgkeep & Renders every text as a marker.
A test aid; never enable it on a site. \\
\opt{quarkus.application.name} & per app & \chgkeep & The name this instance reports in
metrics and traces. \\
\opt{build.timestamp} & build & \chgkeep & Stamped by the build and shown on the
\ui{System} overview. \\
\end{settings}

\subsection{Survey}

\begin{settings}{tab:survey}{Survey: Elicit settings}
\opt{accessCode.autoRegister} & \val{false} & \chgset & Accept \emph{any} access code rather
than only one the console issued. The compose file sets \val{true} for evaluation. A site
delivering real surveys must leave it \val{false}. \\
\opt{elicit.etl.enabled} & \val{true} & \chgmay & Whether this instance builds the reporting
star schema of every installed survey at startup and updates it when a respondent finalizes.
Set \val{false} only on an instance that merely renders surveys and that nothing reports
on. \\
\opt{elicit.etl.drop.enabled} & \val{false} & \chgmay & Whether \val{DELETE
/api/etl/schema/<key>} may drop a survey's reporting schema. Leave it \val{false} at a site:
the next build would regenerate the schema with new surrogate ids, which breaks any extract
keyed on them. Only a database that holds no reporting anyone depends on turns it on. \\
\opt{elicit.diagnostics.startup.enabled} & \val{true} & \chgmay & Whether the startup
diagnostics report is written to the log. \\
\opt{elicit.diagnostics.startup.delay} & \val{10s} & \chgmay & How long after startup that
report is written. \\
\opt{elicit.content.translations-ttl} & \val{PT5M} & \chgmay & How long a survey's content
translations are held in memory before being reloaded, as an ISO-8601 duration. A translation
published through \emph{Apply Survey Definition} reaches respondents within this window without
a restart. See Chapter~\ref{sec:i18n}. \\
\opt{elicit.content.serve-stale-translations} & \val{false} & \chgmay & Whether to show a
translation whose base text has been edited since it was made. Left \val{false}, such a string
falls back to the survey's base language, so every language asks the same question; set
\val{true} only where continuity matters more than that. \\
\opt{elicit.post-survey-action.retry.enabled} & \val{true} & \chgmay & Whether a post-survey
action whose call failed is called again on a schedule. See Section~\ref{sec:psa}. \\
\opt{elicit.post-survey-action.retry.interval} & \val{15m} & \chgmay & How often the retry
runs, as a duration (\val{15m}, \val{1h}). \\
\opt{elicit.post-survey-action.retry.max-attempts} & \val{24} & \chgmay & How many attempts a
delivery gets, the first call included, before Survey stops retrying it. \\
\end{settings}

A post-survey action's credential is set per action, under a name that holds the action's key:
\optlit{elicit.post-survey-action."<key>".bearer-token} (Section~\ref{sec:psa}). It is a
secret and has no default.

\subsection{Admin}

\begin{settings}{tab:admin}{Admin: Elicit settings}
\opt{elicit.authorization.mode} & \val{OIDC} & \chgmay & \val{DATABASE} also consults
\val{survey.user_roles} and shows the role-assignment section of the user editor. It does not
disable OIDC. Environment: \opt{ELICIT_AUTHORIZATION}. \\
\opt{elicit.survey.url} & \val{http://survey:8080} & \chgset & Where Survey is, for the
reporting-schema rebuild after a definition is applied. The default is the compose service
name. Environment: \opt{ELICIT_SURVEY_URL}. \\
\opt{elicit.survey.etl-build.enabled} & \val{true} & \chgmay & Whether to make that call at
all. Environment: \opt{ELICIT_SURVEY_ETL_BUILD_ENABLED}. \\
\opt{elicit.mailer.send-timeout-seconds} & \val{30} & \chgmay & How long a single message
send may take before it is abandoned. \\
\opt{elicit.participation.small-cell} & \val{5} & \chgmay & On the Participation page and in its
CSV and PDF, a count from 1 to one below this value is shown as ``< \textit{k}'' (\textit{k} being this value) and the rates built on it
are withheld, so that a narrow filter cannot single out a respondent. \val{1} turns suppression
off. Environment: \opt{ELICIT_PARTICIPATION_SMALL_CELL}. \\
\opt{elicit.participation.stalled-days} & \val{14} & \chgmay & On the Participation page and in its
exports, a respondent in progress who has not opened a page or saved an answer for this many days
is counted as stalled. Environment: \opt{ELICIT_PARTICIPATION_STALLED_DAYS}. \\
\opt{elicit.debug.reveal-tokens} & \val{false} & \chgkeep & Shows OIDC and bearer tokens
unmasked on the diagnostics screen. A development aid. This concerns identity tokens, not
respondent access codes. \\
\opt{admin.manual.resource} & \val{manual/elicit-admin-manual.pdf} & \chgkeep & Where the
administrator's manual sits on the classpath. A build carrying no manual simply does not
offer it. \\
\opt{quarkus.mailer.mock} & --- & \chgkeep & When true, messages are logged instead of sent.
Set by Quarkus itself in the development and test profiles; never set it on a site. \\
\end{settings}

\subsection{The Quarkus settings the images set}

These are not Elicit settings, but the images set them, so your site is overriding a value
rather than a blank when it changes one.

\begin{settings}{tab:quarkus-db}{Datasources: Survey, Admin and FHHS}
\opt{quarkus.datasource.db-kind} & \val{postgresql} & \chgkeep & Elicit runs on PostgreSQL
only. \\
\opt{quarkus.datasource.db-version} & \val{17} & \chgkeep & The version the schema and the
migrations are written against. \\
\opt{quarkus.datasource.jdbc.url} & \val{jdbc:postgresql://db:5432/survey} plus the
parameters below & \chgset & The application connection. The host names the compose
service. \\
\opt{quarkus.datasource.username} & per app & \chgset & \val{survey_user} for Survey and
FHHS, \val{surveyadmin_user} for Admin. \\
\opt{quarkus.datasource.password} & none & \chgset & No value is shipped. Environment:
\opt{ELICIT_DB_PASSWORD}. See Chapter~\ref{sec:secrets}. \\
\opt{quarkus.datasource.jdbc.min-size} & \val{5} & \chgmay & \\
\opt{quarkus.datasource.jdbc.max-size} & \val{20} & \chgmay & \\
\opt{quarkus.datasource.jdbc.leak-detection-interval} & \val{10m} & \chgmay & \\
\opt{quarkus.datasource.owner.jdbc.url} & as above & \chgset & The migrating connection, as
which every module's Flyway migrations run. \\
\opt{quarkus.datasource.owner.username} & \val{elicit_owner} & \chgkeep & The schema owner. \\
\opt{quarkus.datasource.owner.password} & none & \chgset & No value is shipped.
Environment: \opt{ELICIT_OWNER_DB_PASSWORD}. See Chapter~\ref{sec:secrets}. \\
\end{settings}

The shipped connection URLs carry four parameters that the host and database name above
leave out. Keep them when you point a module at your own server: they are what stops a
dropped connection hanging a request indefinitely:

\begin{elicitfile}[fontsize=\footnotesize]
?connectTimeout=30&socketTimeout=600&tcpKeepAlive=true&tcpKeepAlivesIdle=60
\end{elicitfile}

\begin{settings}{tab:quarkus-oidc}{Identity: Admin only}
\opt{quarkus.oidc.auth-server-url} & \val{http://host.docker.internal:8180/realms/elicit/} &
\chgset & Your provider's issuer address. It must resolve to the same provider from the
browser and from inside the container. \\
\opt{quarkus.oidc.client-id} & \val{elicit-admin} & \chgmay & \\
\opt{quarkus.oidc.credentials.secret} & \opt{OIDC_CLIENT_SECRET} & \chgset & The
client secret. \\
\opt{quarkus.oidc.application-type} & \val{hybrid} & \chgkeep & Accepts browser sign-ins and
bearer tokens on one client. \\
\opt{quarkus.oidc.authentication.redirect-path} & \val{/callback} & \chgkeep & \\
\opt{quarkus.oidc.logout.path} & \val{/logout} & \chgkeep & \\
\end{settings}

\begin{settings}{tab:quarkus-http}{HTTP, proxying and limits}
\opt{quarkus.http.port} & \val{8080} & \chgkeep & Inside the container. Publish a different
host port instead. \\
\opt{quarkus.http.proxy.proxy-address-forwarding} & \val{true} & \chgkeep & Required behind a
load balancer so that redirect URIs are built with the external address. \\
\opt{quarkus.http.proxy.allow-x-forwarded} & \val{true} & \chgkeep & \\
\opt{quarkus.http.cors.origins} & empty & \chgmay & Admin lists allowed origins explicitly;
Survey defaults to \val{http://localhost:8080}. Each application is served same-origin to
its own interface, so a site normally needs nothing here. \\
\opt{quarkus.http.limits.max-body-size} & \val{25M} & \chgmay & Admin only. Bounds a CSV or
definition-file upload. \\
\opt{quarkus.tls.trust-all} & \val{false} & \chgkeep & Survey and Admin verify the
certificate of every service they call over TLS. See the note below. \\
\end{settings}

\begin{elicitnote}[Services signed by a private certificate authority]
Survey and Admin verify every outbound TLS connection (to the identity provider, the mail
relay, report services and each other) against the Java runtime's trust store. Releases before 3.0.0
shipped \opt{quarkus.tls.trust-all}\val{=true}, which turned that verification off; a site
that carried the setting forward should remove it. A service whose certificate is issued by a
private authority is trusted by adding that authority, for example
\opt{quarkus.tls.trust-store.pem.certs} naming its PEM file, not by turning verification off.
\end{elicitnote}

\begin{settings}{tab:quarkus-flyway}{Migrations: Survey and Admin}
\opt{quarkus.flyway.owner.schemas} & \val{survey} & \chgkeep & \\
\opt{quarkus.flyway.owner.migrate-at-start} & per app & \chgkeep & \val{true} on Admin.
\val{false} on Survey, which runs its migrations itself so that it can choose between the
current schema and the upgrade track. \\
\opt{quarkus.flyway.owner.baseline-on-migrate} & \val{true} & \chgkeep & \\
\opt{quarkus.flyway.owner.validate-on-migrate} & \val{true} & \chgkeep & \\
\opt{quarkus.flyway.owner.repair-at-start} & \val{true} & \chgkeep & Admin only. Rewrites the
recorded checksums of two migrations that were edited after release, so existing databases
start without a manual repair. \\
\opt{quarkus.flyway.owner.table} & per app & \chgkeep & \val{flyway_history} for Survey,
\val{flyway_admin_history} for Admin: each module keeps its own history in the shared
schema. \\
\opt{quarkus.flyway.owner.connect-retries} & \val{10} & \chgmay & How many times to retry a
database that is not up yet. \\
\opt{quarkus.flyway.owner.locations} & \val{db/migration} & \chgkeep & Which migration scripts
to run. Survey reads this itself to choose between the current schema and the upgrade track;
changing it points a release at scripts it was not built with. \\
\opt{quarkus.flyway.owner.baseline-version} & \val{0} & \chgkeep & The version recorded when
an existing database is brought under Flyway's control. \\
\opt{quarkus.flyway.owner.baseline-description} & \val{Empty Database} & \chgkeep & The
description recorded alongside it. \\
\end{settings}

\begin{settings}{tab:quarkus-mail}{Mail: Admin only}
\opt{quarkus.mailer.host} & \val{localhost} & \chgset & Your SMTP relay. \\
\opt{quarkus.mailer.port} & \val{25} & \chgset & \\
\opt{quarkus.mailer.from} & --- & \chgset & The address invitations are sent from. \\
\opt{quarkus.mailer.tls} & \val{false} & \chgset & \\
\opt{quarkus.mailer.username} & --- & \chgmay & With
\opt{quarkus.mailer.password} if your relay authenticates. \\
\end{settings}

\begin{settings}{tab:quarkus-obs}{Health, metrics and traces}
\opt{quarkus.smallrye-health.root-path} & \val{/q/health} & \chgkeep & Gives
\opt{/q/health/ready} and \opt{/q/health/live}. Public on Admin. \\
\opt{quarkus.micrometer.export.prometheus.path} & \val{/q/metrics} & \chgkeep & Public on
Admin, for a scraper your site provides. \\
\opt{quarkus.otel.enabled} & \val{true} & \chgmay & \\
\opt{quarkus.otel.exporter.otlp.endpoint} & \val{http://localhost:4317} & \chgset & Your
collector. The compose file sets \val{http://jaeger:4317}. \\
\opt{quarkus.otel.traces.sampler} & \val{parentbased_always_on} & \chgmay & Samples every
trace. Consider reducing this on a busy site. \\
\opt{quarkus.log.level} & \val{WARN} & \chgmay & Environment: \opt{LOG_LEVEL}. \\
\opt{quarkus.log.category."com.elicitsoftware".level} & \val{INFO} & \chgmay & Admin only.
Environment: \opt{ELICIT_LOG_LEVEL}. \\
\opt{quarkus.http.access-log.enabled} & \val{true} & \chgmay & \\
\opt{quarkus.hibernate-orm.statistics} & \val{true} & \chgmay & Useful, and not free. \\
\end{settings}

\subsection{Everything else}

These tables cover what Elicit defines and what its images set. Every other Quarkus setting
behaves as Quarkus documents it:

\begin{center}
\url{https://quarkus.io/guides/all-config}
\end{center}

\end{document}
